\chapter[Dystrybucje i słabe pochodne]{Dystrybucje, słabe pochodne i rozwiązania fundamentalne}
\label{chap:pde-distributions}

W poprzednich rozdziałach rozwiązanie klasyczne musiało mieć tyle pochodnych,
ile wymagało równanie. W praktyce taki warunek jest często zbyt silny:
profil fali może mieć skok, dane początkowe mogą być tylko całkowalne,
a rozwiązanie eliptyczne może być najpierw otrzymane przez minimalizację
energii. Potrzebujemy rachunku różniczkowego, w którym można różniczkować
nawet funkcję skokową. Nie chodzi o arbitralne przypisywanie jej wartości
pochodnej w punkcie skoku, lecz o zastąpienie równań punktowych przez
tożsamości całkowe z funkcjami testowymi.

\section{Przestrzeń funkcji testowych}
\label{sec:pde7-tests}

\begin{definicja}[Funkcje testowe i ich zbieżność]
\label{def:pde7-test}
Niech $d\ge1$ oraz $\Omega\subset\R^d$ będzie niepustym zbiorem otwartym.
Przez $\mathcal D(\Omega)=C_c^\infty(\Omega;\C)$ oznaczamy przestrzeń
zespolonych funkcji nieskończenie różniczkowalnych o zwartym nośniku
zawartym w $\Omega$. Nośnikiem funkcji $\varphi$ jest domknięcie
w $\Omega$ zbioru $\{x\in\Omega:\varphi(x)\ne0\}$.
Notacja $K\Subset\Omega$ oznacza, że $K$ jest zwarty i $K\subset\Omega$.
Dla niepustego $K\Subset\Omega$ definiujemy
\[
 \mathcal D_K(\Omega)=
 \{\varphi\in\mathcal D(\Omega):\supp\varphi\subset K\}
\]
oraz półnormy
\[
 p_{K,m}(\varphi)=\max_{|\alpha|\le m}
 \sup_{x\in K}|\partial^\alpha\varphi(x)|,\qquad m\in\N_0.
\]
Tutaj $\alpha=(\alpha_1,\ldots,\alpha_d)\in\N_0^d$,
$|\alpha|=\sum_{i=1}^d\alpha_i$ oraz
$\partial^\alpha=\partial_{x_1}^{\alpha_1}\cdots
\partial_{x_d}^{\alpha_d}$.
Dla $K=\varnothing$ przyjmujemy $\mathcal D_K(\Omega)=\{0\}$
i półnormy równe zeru.

Przestrzeń $\mathcal D_K(\Omega)$ wyposażamy w topologię
wyznaczoną przez półnormy $p_{K,m}$. Na $\mathcal D(\Omega)$
przyjmujemy standardową topologię lokalnie wypukłą granicy
induktywnej: jest to najsilniejsza topologia lokalnie wypukła,
dla której wszystkie włożenia $\mathcal D_K(\Omega)\to
\mathcal D(\Omega)$ są ciągłe. Wyznaczają ją wszystkie półnormy
na $\mathcal D(\Omega)$, których ograniczenia do każdej
$\mathcal D_K(\Omega)$ są ciągłe. Istotnie, topologia
wyznaczona przez tę rodzinę jest lokalnie wypukła,
a każde włożenie jest ciągłe z definicji rodziny.
W każdej innej topologii lokalnie wypukłej o tej
własności półnormy ją wyznaczające mają ciągłe
ograniczenia do wszystkich $\mathcal D_K(\Omega)$,
więc należą do rozważanej rodziny. Otrzymana
topologia jest zatem najsilniejszą taką topologią.

Ciąg $\varphi_j$ zbiega do $\varphi$ w $\mathcal D(\Omega)$
wtedy i tylko wtedy, gdy istnieje jeden zbiór zwarty
$K\Subset\Omega$ zawierający nośnik $\varphi$ oraz nośniki
wszystkich $\varphi_j$ od pewnego indeksu, a ponadto
$p_{K,m}(\varphi_j-\varphi)\to0$ dla każdego $m\in\N_0$.
\end{definicja}

Uzasadnijmy podany opis zbieżności. Wystarczy rozpatrzyć
$\varphi=0$, ponieważ odejmowanie ustalonej funkcji jest
ciągłe, a suma dwóch zbiorów zwartych jest zwarta.
Jeśli wszystkie rozważane nośniki od pewnego indeksu
zawierają się w $K$ i półnormy dążą do zera, to ciąg
zbiega w $\mathcal D_K(\Omega)$, a ciągłość włożenia
daje zbieżność w $\mathcal D(\Omega)$.

Odwrotnie, niech $\varphi_j\to0$ w topologii granicy
induktywnej. Wybieramy rosnący ciąg zbiorów zwartych
$K_\ell\Subset\Omega$, których wnętrza pokrywają $\Omega$
i które zawierają ostatecznie każdy zbiór zwarty w $\Omega$.
Można wziąć
\[
 K_\ell=\{x\in\Omega:|x|\le\ell,\
 \operatorname{dist}(x,\R^d\setminus\Omega)\ge1/\ell\},
 \qquad \ell\ge1,
\]
przyjmując odległość od zbioru pustego równą $+\infty$.
Gdy dopełnienie jest niepuste, jego funkcja odległości
jest ciągła na $\R^d$, więc $K_\ell$ jest domknięty
i ograniczony oraz zawarty w $\Omega$.
Dla $\Omega=\R^d$ otrzymujemy domknięte kule.
Każdy niepusty zwarty podzbiór $\Omega$ jest ograniczony
i ma dodatnią odległość od niepustego dopełnienia
$\Omega$: ciągła funkcja odległości osiąga na nim
dodatnie minimum. Wystarczy więc wybrać $\ell$
z oboma warunkami definiującymi $K_\ell$ spełnionymi
ostro na tym podzbiorze. W szczególności wnętrza
$K_\ell$ pokrywają $\Omega$.
Jeśli nie istnieje jeden wspólny zwarty zbiór dla nośników
od pewnego indeksu, możemy kolejno wybrać rosnące indeksy
$j_\ell$ i punkty $x_\ell\notin K_\ell$ takie, że
$\varphi_{j_\ell}(x_\ell)\ne0$. Istnienie punktu wynika
z tego, że $K_\ell$ jest domknięty i nie zawiera
nośnika $\varphi_{j_\ell}$.
Definiujemy półnormę
\[
 q(\psi)=\sum_{\ell\ge1}
 \frac{|\psi(x_\ell)|}{|\varphi_{j_\ell}(x_\ell)|},
 \qquad \psi\in\mathcal D(\Omega).
\]
Dla każdej ustalonej $\psi$ suma ma tylko skończenie
wiele niezerowych składników: jej nośnik jest zawarty
w pewnym $K_{\ell_0}$, a dla $\ell\ge\ell_0$
punkt $x_\ell$ leży poza tym zbiorem. Na każdej
$\mathcal D_K(\Omega)$ półnorma $q$ jest zatem
ograniczona przez pewną skończoną wielokrotność $p_{K,0}$.
Z definicji topologii granicy induktywnej $q$ jest ciągła.
Jednak $q(\varphi_{j_\ell})\ge1$ dla każdego $\ell$,
co przeczy zbieżności do zera. Istnieje więc wspólny
zwarty zbiór $K$ zawierający nośniki od pewnego indeksu.

Dla każdego $m\in\N_0$ półnorma
\[
 q_m(\psi)=\max_{|\alpha|\le m}
 \sup_{x\in\Omega}|\partial^\alpha\psi(x)|
\]
jest skończona i ciągła w topologii granicy induktywnej:
na $\mathcal D_K(\Omega)$ pokrywa się z $p_{K,m}$,
ponieważ wszystkie pochodne znikają poza nośnikiem
funkcji. Zatem $q_m(\varphi_j)\to0$, a na ustalonym
$K$ otrzymujemy wymagane $p_{K,m}(\varphi_j)\to0$.
Dowodzi to obu implikacji.

Warto odróżnić zwarty nośnik w $\Omega$ od samej
gładkości na domknięciu $\Omega$. Jeśli $\Omega\ne\R^d$,
zwartość nośnika i jego rozłączność z domkniętym
dopełnieniem $\Omega$ zapewniają dodatnią odległość
od tego dopełnienia (dla funkcji zerowej wniosek jest
bezpośredni). Funkcja testowa znika więc w pewnym
otoczeniu brzegu od strony $\Omega$, a przedłużenie
jej przez zero na $\R^d$ jest gładkie. Gdy
$\Omega=\R^d$, przedłużanie nie jest potrzebne.

Istnieje wiele funkcji testowych. Dla przykładu funkcja
\[
 \eta(x)=\begin{cases}
 \exp\!\bigl(-1/(1-|x|^2)\bigr),& |x|<1,\\
 0,& |x|\ge1
 \end{cases}
\]
jest klasy $C_c^\infty(\R^d)$ i jest dodatnia w $B_1(0)$.
Aby to uzasadnić, połóżmy
\[
 \rho(s)=
 \begin{cases}e^{-1/s},&s>0,\\0,&s\le0.\end{cases}
\]
Dla $s>0$ każda pochodna $\rho^{(m)}(s)$ ma postać
$P_m(1/s)e^{-1/s}$, gdzie $P_m$ jest wielomianem:
dla $m=0$ mamy $P_0=1$, a różniczkowanie przy
$z=1/s$ daje $P_{m+1}(z)=z^2(P_m(z)-P_m'(z))$.
Dla dowolnego całkowitego $r\ge0$ zachodzi
$z^r e^{-z}\to0$ przy $z\to+\infty$, na przykład
z oszacowania $e^z\ge z^{r+1}/(r+1)!$.
Każda z powyższych pochodnych, a także jej iloraz
przez $s$, dąży więc do zera przy $s\downarrow0$.
Indukcyjnie ilorazy różnicowe w zerze pokazują,
że $\rho$ jest klasy $C^\infty(\R)$ i wszystkie
jej pochodne w zerze znikają.
Ponieważ $\eta(x)=\rho(1-|x|^2)$, gładkość wynika
z reguły różniczkowania złożenia. Nośnikiem
$\eta$ jest domknięta kula jednostkowa.
Jeśli $\overline{B_r(a)}\subset\Omega$, $r>0$,
to funkcja $x\mapsto\eta((x-a)/r)$ jest funkcją
testową w $\Omega$, dodatnią na $B_r(a)$.

\begin{uwaga}[Dlaczego potrzebny jest wspólny nośnik?]
Dla $d=1$ przesunięcia $\varphi_j(x)=\eta(x-j)$ na $\R$
mają wszystkie pochodne jednostajnie ograniczone
względem $j$, a na każdym ustalonym zbiorze zwartym
funkcje i ich pochodne od pewnego indeksu znikają.
Mimo to nie dążą do zera w $\mathcal D(\R)$, ponieważ
$\supp\varphi_j=[j-1,j+1]$ i żaden zbiór zwarty
nie zawiera tych nośników od pewnego indeksu.
Sama zbieżność jednostajna wszystkich pochodnych
na zbiorach zwartych nie jest zatem zbieżnością
w przestrzeni funkcji testowych.
\end{uwaga}

\section{Dystrybucje jako ciągłe funkcjonały liniowe}
\label{sec:pde7-def}

\begin{definicja}[Dystrybucja]
\label{def:pde7-distribution}
Dla niepustego otwartego $\Omega\subset\R^d$ dystrybucją na $\Omega$
nazywamy ciągły funkcjonał \emph{liniowy} $T:\mathcal D(\Omega)\to\C$.
Zbiór wszystkich takich funkcjonałów oznaczamy przez $\mathcal D'(\Omega)$.
Wartość funkcjonału zapisujemy jako $\langle T,\varphi\rangle$;
nie jest to na ogół iloczyn skalarny, a w szczególności
zapis jest liniowy względem $\varphi$, bez sprzężenia zespolonego.
\end{definicja}

Związek z definicją topologiczną można wyrazić za pomocą lokalnych
oszacowań. Jest to pierwszy ważny przykład tego, że element
przestrzeni sprzężonej do przestrzeni funkcji można opisywać
przez oszacowania na jej półnormach.

\begin{twierdzenie}[Lokalne oszacowanie dystrybucji]
\label{thm:pde7-local-order}
Niech $\Omega\subset\R^d$ będzie otwarta, a $T$ liniowym funkcjonałem
na $\mathcal D(\Omega)$. Wówczas $T\in\mathcal D'(\Omega)$ wtedy i tylko
wtedy, gdy dla każdego $K\Subset\Omega$ istnieją $m_K\in\N_0$
oraz $C_K>0$ takie, że
\begin{equation}\label{eq:pde7-estimate}
 |\langle T,\varphi\rangle|\le C_K p_{K,m_K}(\varphi),
 \qquad \varphi\in\mathcal D_K(\Omega).
\end{equation}
Dla $K\ne\varnothing$ prawa strona oznacza
$C_K\max_{|\alpha|\le m_K}\sup_{x\in K}
|\partial^\alpha\varphi(x)|$, a dla $K=\varnothing$
jest równa zeru zgodnie z definicją półnorm.
\end{twierdzenie}
\begin{proof}
Jeżeli $T$ jest ciągły, to przez ciągłość włożenia
jego ograniczenie do $\mathcal D_K(\Omega)$ jest ciągłe.
Ponieważ półnormy $p_{K,m}$ rosną względem $m$,
każde otoczenie zera zawiera zbiór postaci
$\{\varphi:p_{K,m}(\varphi)<\delta\}$ dla pewnych
$m$ i $\delta>0$ (w skończonym przecięciu warunków
wybieramy największy indeks i najmniejszy promień).
Ciągłość w zerze daje zatem $|T(\varphi)|<1$,
gdy $p_{K,m}(\varphi)<\delta$.
Jeżeli $p_{K,m}(\varphi)>0$, stosujemy tę implikację
do funkcji $\delta\varphi/(2p_{K,m}(\varphi))$.
Z liniowości otrzymujemy
$|T(\varphi)|\le(2/\delta)p_{K,m}(\varphi)$.
Jeżeli $p_{K,m}(\varphi)=0$, to w szczególności
$p_{K,0}(\varphi)=0$; funkcja znika zarówno na $K$,
jak i poza $K$, więc jest zerowa. Dla pustego $K$
oszacowanie również zachodzi, gdyż przestrzeń
$\mathcal D_K(\Omega)$ jest trywialna.

Odwrotnie, jeśli zachodzi~\eqref{eq:pde7-estimate},
to półnorma $\varphi\mapsto|T(\varphi)|$ ma ciągłe
ograniczenie do każdej $\mathcal D_K(\Omega)$.
Z charakterystyki topologii granicy induktywnej
w sekcji~\ref{sec:pde7-tests} jest ona ciągła na
$\mathcal D(\Omega)$. To daje ciągłość $T$.

Równoważnie, funkcjonał liniowy jest dystrybucją,
jeśli zachowuje zbieżność ciągów do zera. Istotnie,
ciągłość implikuje tę własność. Gdyby dla pewnego
$K$ nie zachodziło~\eqref{eq:pde7-estimate}, to dla
każdego $j\ge1$ istniałaby $\psi_j\in\mathcal D_K(\Omega)$
taka, że $|T(\psi_j)|>j p_{K,j}(\psi_j)$.
Wtedy $\varphi_j=\psi_j/T(\psi_j)$ spełnia
$T(\varphi_j)=1$ i $p_{K,j}(\varphi_j)<1/j$.
Dla każdego ustalonego $m$ i $j\ge m$ mamy
$p_{K,m}(\varphi_j)\le p_{K,j}(\varphi_j)$, więc
$\varphi_j\to0$ w $\mathcal D(\Omega)$, sprzecznie
z zachowaniem zbieżności do zera.
\end{proof}

Liczba $m_K$ może zależeć od $K$. Mówimy wtedy o \emph{lokalnie
skończonym rzędzie} dystrybucji. Nie należy bez dodatkowych założeń
wnioskować, że jedna liczba $m$ działa na wszystkich zwartych
podzbiorach obszaru, \emph{nawet gdy jest on ograniczony}.
Istotnie, na $\Omega=(0,1)$ określmy funkcjonał
\[
 \langle T,\varphi\rangle
 =\sum_{k=2}^\infty 2^{-k^2}(-1)^k\varphi^{(k)}(1/k).
\]
W następnej sekcji zapiszemy go jako
$\sum_{k=2}^\infty 2^{-k^2}\partial^k\delta_{1/k}$.
Dla każdego $K\Subset(0,1)$ zbiór
$I_K=\{k\ge2:1/k\in K\}$ jest skończony.
Jeśli $\supp\varphi\subset K$, to wszystkie pochodne
$\varphi$ znikają poza $K$, więc suma zawiera tylko
składniki o indeksach z $I_K$. Jest zatem liniowa
i dla niepustego $I_K$ spełnia oszacowanie
\[
 |T(\varphi)|\le
 \left(\sum_{k\in I_K}2^{-k^2}\right)
 p_{K,\max I_K}(\varphi).
\]
Dla pustego $I_K$ funkcjonał na $\mathcal D_K$ znika.
Z twierdzenia~\ref{thm:pde7-local-order} wynika, że
$T$ jest dystrybucją.

Nie istnieje jednak wspólny skończony rząd $m$.
Ustalmy $k\ge2$, oznaczmy $a=1/k$ i wybierzmy
$r>0$ tak, aby $K=[a-r,a+r]\subset(0,1)$
nie zawierał żadnego innego punktu $1/\ell$.
Wybierzmy $\theta\in C_c^\infty((-1,1))$ z $\theta(0)=1$
(np. $\theta(s)=\eta(2s)/\eta(0)$
z poprzedniej sekcji w wymiarze jeden). Niech
$\psi(s)=s^k\theta(s)/k!$.
Wtedy $\psi^{(k)}(0)=1$ z reguły Leibniza.
Dla $0<\varepsilon<\min\{1,r\}$ połóżmy
\[
 \varphi_\varepsilon(x)=\varepsilon^k
 \psi((x-a)/\varepsilon).
\]
Jej nośnik zawiera się w $K$, a dla każdego $j\ge0$
mamy $\varphi_\varepsilon^{(j)}(x)=
\varepsilon^{k-j}\psi^{(j)}((x-a)/\varepsilon)$.
Stąd dla $m<k$ zachodzi
$p_{K,m}(\varphi_\varepsilon)\to0$ przy
$\varepsilon\downarrow0$, podczas gdy
\[
 \langle T,\varphi_\varepsilon\rangle
 =2^{-k^2}(-1)^k
\]
nie zależy od $\varepsilon$ i nie jest zerem.
Na tym $K$ najmniejszy dopuszczalny rząd wynosi
dokładnie $k$: oszacowanie rzędu $k$ wynika
bezpośrednio z jedynego składnika sumy, a niższe
rzędy wyklucza skonstruowany ciąg (np. $\varepsilon\to0$
wzdłuż dowolnego ciągu dodatniego).
Ponieważ $k$ jest dowolne, ograniczoność zbioru
otwartego nie zapewnia globalnego skończonego rzędu.

\begin{przyklad}[Dystrybucja regularna i miara punktowa]
\label{ex:pde7-regular}
Jeśli $f\in L^1_{\mathrm{loc}}(\Omega)$, to wzór (niezależny od wyboru reprezentanta klasy funkcji $f$)
\begin{equation}\label{eq:pde7-regular}
 \langle T_f,\varphi\rangle
 =\int_\Omega f(x)\varphi(x)\,\mathrm d\lambda_d(x)
\end{equation}
definiuje dystrybucję, ponieważ dla $\supp\varphi\subset K$ zachodzi
$|T_f(\varphi)|\le\|\varphi\|_\infty\int_K|f(x)|\,\mathrm d\lambda_d(x)$.
Oba warunki: liniowość i ciągłość, są więc spełnione.
Dla $a\in\Omega$ wzór $\langle\delta_a,\varphi\rangle=\varphi(a)$
definiuje \emph{deltę Diraca}. Jest to funkcjonał liniowy:
dla $\supp\varphi\subset K$ mamy
$|\varphi(a)|\le p_{K,0}(\varphi)$, gdy $a\in K$,
a gdy $a\notin K$, wartość jest zerowa.
Zatem delta Diraca ma rząd zerowy.
\end{przyklad}

\begin{stwierdzenie}[Iniektywność odwzorowania $f\mapsto T_f$]
\label{prop:pde7-injective}
Jeżeli $f,g\in L^1_{\mathrm{loc}}(\Omega)$ i $T_f=T_g$, to $f=g$
prawie wszędzie w $\Omega$ względem miary Lebesgue'a.
\end{stwierdzenie}
\begin{proof}
Wystarczy rozważyć $h=f-g$. Wybierzmy $a\in\Omega$
i $r>0$ tak, aby $\overline{B_{2r}(a)}\Subset\Omega$.
Ustalmy nieujemną funkcję $\rho\in C_c^\infty(B_1(0))$
o całce $1$ (skalujemy funkcję $\eta$ do kuli
o promieniu $1/2$ i dzielimy przez jej dodatnią,
skończoną całkę). Połóżmy
$\rho_\eps(x)=\eps^{-d}\rho(x/\eps)$ dla $0<\eps<r$.
Dla każdego $x\in B_r(a)$ funkcja
$y\mapsto\rho_\eps(x-y)$ jest testowa i ma nośnik
w $B_{2r}(a)$. Ponieważ $T_h=0$, mamy
$(h*\rho_\eps)(x)=0$ dla wszystkich takich $x$.
Aby przejść do granicy, przedłużmy $h$ po obcięciu do
$B_{2r}(a)$ przez zero i oznaczmy otrzymaną funkcję
przez $H\in L^1(\R^d)$. Z nierówności Minkowskiego
wynika
\[
 \|H*\rho_\eps-H\|_{L^1(\R^d)}
 \le\int_{B_1(0)}\rho(z)
 \|H(\cdot-\eps z)-H\|_{L^1(\R^d)}
 \,\mathrm d\lambda_d(z)\longrightarrow0.
\]
Istotnie, dla $v\in C_c(\R^d)$ i $|b|\le1$
funkcja $v(\cdot-b)-v$ ma nośnik w ustalonym
zbiorze zwartym (suma nośnika $v$ i domkniętej
kuli jednostkowej). Jednostajna ciągłość $v$
daje więc $\|v(\cdot-b)-v\|_1\to0$ dla $b\to0$.
Korzystamy z gęstości $C_c(\R^d)$ w $L^1(\R^d)$.
Uzasadnijmy ją bez odwoływania się do wygładzania.
Funkcje proste o nośniku skończonej miary są gęste
w $L^1$: obcinamy funkcję do ograniczonej kuli,
następnie ograniczamy jej moduł i aproksymujemy
część rzeczywistą oraz urojoną funkcjami prostymi
z dowolnie małym błędem jednostajnym na tej kuli.
Dwa pierwsze błędy w normie $L^1$ dążą do zera
z całkowalności funkcji, a ostatni jest ograniczony
przez błąd jednostajny razy miarę kuli.
Dla mierzalnego zbioru $A$ skończonej miary
i $\tau>0$ regularność miary Lebesgue'a daje
zwarty $K_0\subset A$ i otwarty $O\supset A$
z $\lambda_d(O\setminus K_0)<\tau$.
Jeśli $K_0=\varnothing$, funkcja zerowa
aproksymuje $\chi_A$ z błędem mniejszym niż $\tau$.
W przeciwnym razie wybieramy $r_0>0$ mniejszy
od odległości $K_0$ od $\R^d\setminus O$
(dowolny skończony $r_0>0$, gdy $O=\R^d$).
Funkcja
\[
 w(x)=\max\{0,1-\operatorname{dist}(x,K_0)/r_0\}
\]
jest ciągła, równa $1$ na $K_0$, ma zwarty
nośnik w $O$ i przyjmuje wartości w $[0,1]$.
Zatem $|\chi_A-w|\le\chi_{O\setminus K_0}$,
skąd $\|\chi_A-w\|_1<\tau$.
Stosując tę konstrukcję do składników funkcji
prostej i nierówność trójkąta, uzyskujemy gęstość. Dla dowolnego $\gamma>0$
wybieramy $v$ z $\|H-v\|_1<\gamma$. Niezmienniczość
miary Lebesgue'a względem przesunięć daje
\[
 \|H(\cdot-b)-H\|_1
 \le2\gamma+\|v(\cdot-b)-v\|_1.
\]
Z dowolności $\gamma$ wynika ciągłość przesunięć.
W oszacowaniu splotu stosujemy twierdzenie
o zbieżności zmajoryzowanej z całkowalną
majorantą $2\|H\|_1\rho(z)$.
Samo oszacowanie wynika z podstawienia
$y=x-\eps z$, normalizacji $\rho$ i twierdzenia
Tonellego zastosowanego do modułu różnicy.

Na $B_r(a)$ sploty $H*\rho_\eps$ są zerowe,
ponieważ wszystkie używane punkty $y$ leżą
w $B_{2r}(a)$, gdzie $H=h$ prawie wszędzie.
Dlatego $\|h\|_{L^1(B_r(a))}
\le\|H*\rho_\eps-H\|_{L^1(\R^d)}\to0$.
Zatem $h=0$ prawie wszędzie na tej kuli.
Z drugiej przeliczalności $\R^d$ wybieramy
przeliczalne pokrycie $\Omega$ takimi kulami.
Przeliczalna suma zbiorów miary zero ma miarę zero,
co daje tezę na całym $\Omega$.
Ogólną teorię przybliżonych jedynek rozwiniemy później
w sekcji~\ref{sec:pde7-mollifiers}; nie jest ona
przesłanką niniejszego dowodu.
\end{proof}

\begin{uwaga}[Dlaczego delta nie jest zwykłą funkcją?]
Załóżmy, że $\delta_0=T_f$ dla pewnego
$f\in L^1_{\mathrm{loc}}(\R^d)$. Wybierzmy funkcję
$\psi\in\mathcal D(\R^d)$, $0\le\psi\le1$,
$\psi(0)=1$, $\supp\psi\subset B_1(0)$
(na przykład $\psi(x)=\eta(2x)/\eta(0)$), i połóżmy
$\psi_\varepsilon(x)=\psi(x/\varepsilon)$.
Wtedy $\langle\delta_0,\psi_\varepsilon\rangle=1$,
natomiast dla $0<\varepsilon<1$
\[
 |\langle T_f,\psi_\varepsilon\rangle|
 \le\int_{B_\varepsilon(0)}|f(x)|\,\mathrm d\lambda_d(x)
 \longrightarrow0.
\]
Istotnie, $f$ jest całkowalna na domkniętej kuli
jednostkowej, a miara $B_\varepsilon(0)$ dąży do zera,
więc stosujemy absolutną ciągłość całki Lebesgue'a.
Otrzymujemy sprzeczność.
\end{uwaga}

\section{Pochodne dystrybucyjne i skoki}
\label{sec:pde7-derivatives}

\begin{definicja}[Pochodna dystrybucyjna]
\label{def:pde7-derivative}
Dla $T\in\mathcal D'(\Omega)$ i wielowskaźnika
$\alpha\in\N_0^d$ określamy $\partial^\alpha T\in\mathcal D'(\Omega)$
wzorem
\begin{equation}\label{eq:pde7-distributional-derivative}
 \langle\partial^\alpha T,\varphi\rangle
 =(-1)^{|\alpha|}\langle T,\partial^\alpha\varphi\rangle,
 \qquad\varphi\in\mathcal D(\Omega).
\end{equation}
\end{definicja}

Definicja ma sens: nośnik $\partial^\alpha\varphi$
zawiera się w nośniku $\varphi$, a dla $\varphi\in\mathcal D_K(\Omega)$
mamy
\[
 p_{K,m}(\partial^\alpha\varphi)
 \le p_{K,m+|\alpha|}(\varphi).
\]
Zatem jeśli $T$ spełnia lokalne oszacowanie
~\eqref{eq:pde7-estimate} z $m_K,C_K$, to
\[
 |\langle\partial^\alpha T,\varphi\rangle|
 \le C_K p_{K,m_K+|\alpha|}(\varphi).
\]
Funkcjonał $\partial^\alpha T$ jest liniowy i ciągły
na mocy twierdzenia~\ref{thm:pde7-local-order}.
W szczególności każda dystrybucja ma pochodne
dowolnego rzędu. Dla wielowskaźników $\alpha,\beta$
dwukrotne zastosowanie definicji daje
\[
 \langle\partial^\beta\partial^\alpha T,\varphi\rangle
 =(-1)^{|\alpha|+|\beta|}
 \langle T,\partial^\alpha\partial^\beta\varphi\rangle
 =\langle\partial^{\alpha+\beta}T,\varphi\rangle.
\]
W ostatnim kroku wykorzystaliśmy komutowanie
pochodnych gładkiej funkcji testowej. Pochodne
mieszane dystrybucji zatem również komutują.

\begin{stwierdzenie}[Zgodność z klasycznym różniczkowaniem]
\label{prop:pde7-classical}
Jeśli $u\in C^1(\Omega)$, to $\partial_jT_u=T_{\partial_j u}$.
Jeśli $T\in\mathcal D'(\Omega)$ i $a\in C^\infty(\Omega)$,
to iloczyn $aT$, dany wzorem
$\langle aT,\varphi\rangle=\langle T,a\varphi\rangle$,
jest dystrybucją i zachodzi reguła Leibniza
$\partial_j(aT)=(\partial_ja)T+a\partial_jT$.
\end{stwierdzenie}
\begin{proof}
Dla $u\in C^1(\Omega)$ iloczyn $u\varphi$ ma zwarty
nośnik w $\Omega$ i przedłuża się przez zero do
funkcji klasy $C^1$ na $\R^d$. Z twierdzenia Fubiniego
i podstawowego twierdzenia rachunku całkowego
całka z jej pochodnej względem $x_j$ jest zerowa
(całkujemy najpierw po całej prostej w tej zmiennej).
Reguła Leibniza dla funkcji daje więc
\[
 \langle\partial_jT_u,\varphi\rangle
 =-\int_\Omega u(x)\partial_j\varphi(x)\,\mathrm d\lambda_d(x)
 =\int_\Omega \partial_ju(x)\varphi(x)\,\mathrm d\lambda_d(x).
\]
To dowodzi pierwszej tezy.

Dla $a\in C^\infty(\Omega)$ iloczyn $a\varphi$ ma
nośnik zawarty w nośniku $\varphi$. Z wielowskaźnikowej
reguły Leibniza wynika dla niepustego $K$
\[
 p_{K,m}(a\varphi)\le 2^m p_{K,m}^{\,a}p_{K,m}(\varphi),
 \qquad
 p_{K,m}^{\,a}=\max_{|\gamma|\le m}
 \sup_{x\in K}|\partial^\gamma a(x)|.
\]
Istotnie, dla każdej pochodnej $\partial^\alpha$
sumę współczynników dwumianowych w regule Leibniza
stanowi $2^{|\alpha|}\le2^m$. Wszystkie pochodne
$a$ są ograniczone na $K$. Stosując lokalne oszacowanie
dla $T$ do $a\varphi$, otrzymujemy ciągłość liniowego
funkcjonału $aT$ (dla pustego $K$ jest to trywialne).
Ponadto
\[
 \begin{aligned}
 \langle(\partial_ja)T+a\partial_jT,\varphi\rangle
 &=\langle T,(\partial_ja)\varphi\rangle
   -\langle T,\partial_j(a\varphi)\rangle\\
 &=-\langle T,a\partial_j\varphi\rangle
 =\langle\partial_j(aT),\varphi\rangle.
 \end{aligned}
\]
Jest to żądana reguła Leibniza, poprawna także
dla zespolonych $a$ i $\varphi$.
\end{proof}

\begin{przyklad}[Funkcja Heaviside'a]
\label{ex:pde7-heaviside}
Niech $H(x)=0$ dla $x<0$ i $H(x)=1$ dla $x>0$;
wartość w zerze nie ma znaczenia. Dla $\varphi\in\mathcal D(\R)$
\[
 \langle H',\varphi\rangle
 =-\int_0^\infty\varphi'(x)\,\mathrm d\lambda_1(x)
 =\varphi(0)=\langle\delta_0,\varphi\rangle.
\]
Zatem $H'=\delta_0$ w $\mathcal D'(\R)$.
Utożsamiamy tu funkcję lokalnie całkowalną z jej
dystrybucją regularną. Dla $u(x)=|x|$ całkujemy osobno
na obu półprostych:
\[
 \begin{aligned}
 \langle T_u',\varphi\rangle
 &=-\int_{-\infty}^0(-x)\varphi'(x)\,\mathrm d\lambda_1(x)
   -\int_0^\infty x\varphi'(x)\,\mathrm d\lambda_1(x)\\
 &=-\int_{-\infty}^0\varphi(x)\,\mathrm d\lambda_1(x)
   +\int_0^\infty\varphi(x)\,\mathrm d\lambda_1(x).
 \end{aligned}
\]
Wyrazy brzegowe w zerze znikają dzięki czynnikowi $x$,
a w nieskończoności dzięki zwartemu nośnikowi.
Zatem $T_u'=T_{\operatorname{sgn}}$. Ponieważ
$\operatorname{sgn}=2H-1$ prawie wszędzie i pochodna
dystrybucji stałej jest zerowa, mamy $T_u''=2\delta_0$.
\end{przyklad}

\begin{przyklad}[Pochodna delty]
Bezpośrednio z definicji pochodnej otrzymujemy
\[
 \langle\partial^\alpha\delta_a,\varphi\rangle
 =(-1)^{|\alpha|}\partial^\alpha\varphi(a).
\]
Jest to dystrybucja rzędu co najwyżej $|\alpha|$,
a na każdym zwartym $K$ zawierającym $a$ w swoim
wnętrzu jej rząd wynosi dokładnie $|\alpha|$.
Aby uzasadnić dolne ograniczenie, dla $k=|\alpha|>0$
wybieramy $\theta\in C_c^\infty(B_1(0))$ z $\theta(0)=1$
i kładziemy $\psi(x)=x^\alpha\theta(x)/\alpha!$,
gdzie $x^\alpha=\prod_{i=1}^d x_i^{\alpha_i}$ oraz
$\alpha!=\prod_{i=1}^d\alpha_i!$.
Reguła Leibniza daje $\partial^\alpha\psi(0)=1$.
Dla dostatecznie małego $\varepsilon>0$ funkcje
$\varphi_\varepsilon(x)=\varepsilon^k
\psi((x-a)/\varepsilon)$ należą do $\mathcal D_K(\Omega)$
i spełniają $p_{K,m}(\varphi_\varepsilon)\to0$
dla $m<k$, natomiast
$\langle\partial^\alpha\delta_a,\varphi_\varepsilon\rangle
=(-1)^k$. Wyklucza to oszacowania niższego rzędu.
Dla $k=0$ otrzymujemy deltę rzędu zerowego.

W szczególności $\delta'_0$ nie jest dystrybucją
wyznaczoną przez lokalnie skończoną miarę Radona
(również zespoloną): taka miara $\mu$ spełnia
$|\int\varphi\,\mathrm d\mu|\le|\mu|(K)p_{K,0}(\varphi)$,
gdzie $|\mu|$ oznacza jej wahanie całkowite na $K$.
Dla $\delta'_0$ właśnie wykazaliśmy niemożliwość
oszacowania rzędu zerowego.
\end{przyklad}

\begin{twierdzenie}[Wzór skoku przez hiperpowierzchnię]
\label{thm:pde7-jump}
Niech $\Omega\subset\R^d$ będzie otwarta, a
$S\subset\Omega$ będzie względnie domkniętą,
zorientowaną, zanurzoną hiperpowierzchnią klasy $C^1$
bez brzegu w $\Omega$. Zakładamy, że w każdym
punkcie $S$ istnieje otoczenie, w którym $S$ jest
wykresem funkcji klasy $C^1$ (po obrocie współrzędnych)
i rozdziela to otoczenie na dwie strony $\Omega_-$
i $\Omega_+$. Nazwy stron są zgodne z orientacją:
jednostkowa normalna $\nu$ jest skierowana
od strony minus do strony plus.
Niech $u\in C^1(\Omega\setminus S)$ i niech lokalnie
przy każdym punkcie $S$ ograniczenia $u$ do obu stron
mają przedłużenia klasy $C^1$ na ich domknięcia.
Oznaczmy ślady na $S$ przez $u_-$ i $u_+$.
Wtedy, w sensie dystrybucji,
\begin{equation}\label{eq:pde7-jump}
 \partial_jT_u=T_{(\partial_ju)_{\mathrm{reg}}}
 +(u_+-u_-)\nu_j\delta_S,
\end{equation}
gdzie $(\partial_ju)_{\mathrm{reg}}$ oznacza klasyczną
pochodną poza $S$, a dla ciągłej funkcji $b$ na $S$
\[
 \langle b\delta_S,\varphi\rangle
 =\int_S b(s)\varphi(s)\,\mathrm d\sigma(s),
 \qquad \sigma=\mathcal H^{d-1}\!\restriction_S.
\]
Wartości funkcji objętościowych na $S$ nie wpływają
na ich dystrybucje regularne.
\end{twierdzenie}
\begin{proof}
W lokalnych współrzędnych wykresowych hiperpowierzchnia
ma miarę Lebesgue'a zero z twierdzenia Fubiniego:
każdy przekrój w kierunku osi wykresu jest co najwyżej
jednopunktowy. Przeliczalne pokrycie takimi otoczeniami
daje $\lambda_d(S)=0$. Na każdym zwartym $K\Subset\Omega$
zbiór $K\cap S$ jest zwarty dzięki względnemu domknięciu
$S$. Można go pokryć skończenie wieloma mniejszymi
otoczeniami wykresowymi o ograniczonej podstawie
i ograniczonej pochodnej funkcji wykresu.
Wzór na miarę powierzchni wykresu
$\sqrt{1+|\nabla h|^2}\,\mathrm d\lambda_{d-1}$
daje $\sigma(K\cap S)<\infty$.
W wymiarze $d=1$ stosujemy zamiast tego miarę
liczącą: $K\cap S$ jest skończonym zbiorem punktów.
Z lokalnych przedłużeń klasy $C^1$ i skończonego
pokrycia $K$ wynika lokalna całkowalność $u$
i $(\partial_ju)_{\mathrm{reg}}$. Ponadto
\[
 |\langle b\delta_S,\varphi\rangle|
 \le \sup_{K\cap S}|b|\,\sigma(K\cap S)p_{K,0}(\varphi)
\]
dla $\supp\varphi\subset K$ (gdy przecięcie jest
puste, całka znika). Zatem wyraz powierzchniowy
jest dystrybucją rzędu zerowego.

Ustalmy $\varphi\in\mathcal D(\Omega)$.
Najpierw załóżmy, że jej nośnik jest zawarty
w jednym otoczeniu wykresowym. Możemy wybrać
lokalny cylinder wykresowy o podstawie prostokątnej,
zawierający nośnik w swoim wnętrzu, którego górna
i dolna ściana nie przecinają wykresu.
Wzór Gaussa na obu częściach cylindra stosujemy
do pola $u\varphi\,e_j$, gdzie $e_j$ jest $j$-tym
wektorem standardowej bazy. Obie części mają
brzeg odcinkami klasy $C^1$, a pole ma ciągłe
pochodne aż do strony wykresu.
Na sztucznych ścianach cylindra pole znika,
ponieważ $\varphi$ ma nośnik we wnętrzu.
Normalna zewnętrzna części minus na $S$ to $\nu$,
a części plus to $-\nu$. W konsekwencji
\[
 \begin{aligned}
 -\int_{\Omega_-}u\partial_j\varphi\,\mathrm d\lambda_d
 &=\int_{\Omega_-}(\partial_ju)\varphi\,\mathrm d\lambda_d
   -\int_S u_-\varphi\nu_j\,\mathrm d\sigma,\\
 -\int_{\Omega_+}u\partial_j\varphi\,\mathrm d\lambda_d
 &=\int_{\Omega_+}(\partial_ju)\varphi\,\mathrm d\lambda_d
   +\int_S u_+\varphi\nu_j\,\mathrm d\sigma.
 \end{aligned}
\]
Tutaj całki dotyczą części lokalnego cylindra.
Sumując, otrzymujemy~\eqref{eq:pde7-jump} dla takiego
testu. Dla testu o nośniku rozłącznym z $S$ wzór
wynika z klasycznego całkowania przez części.

W ogólnym przypadku pokrywamy zwarty nośnik
$\varphi$ skończoną rodziną otoczeń wykresowych
i otoczeń rozłącznych z $S$. Korzystamy z gładkiego
podziału jedności podporządkowanego temu pokryciu
w otoczeniu nośnika. Dokładniej, wybieramy skończenie
wiele nieujemnych funkcji testowych $\eta_i$
o nośnikach w tych otoczeniach, których suma $A$
jest dodatnia w otoczeniu nośnika $\varphi$
(istnienie wynika z funkcji na kulach z sekcji
~\ref{sec:pde7-tests} i zwartości nośnika).
Funkcje $\varphi_i=\varphi\eta_i/A$ tam, gdzie $A>0$,
przedłużamy przez zero poza ten zbiór.
Są gładkie, gdyż nośnik $\varphi$ jest zawarty
w otwartym zbiorze $\{A>0\}$, oraz mają zwarte
nośniki w odpowiednich otoczeniach.
Mamy $\varphi=\sum_i\varphi_i$. Stosując wzór
lokalny do każdej $\varphi_i$ i sumując skończenie
wiele równości, otrzymujemy tezę dla $\varphi$.
\end{proof}

\begin{wniosek}[Skok strumienia i równanie ciągłości]
Niech $F=(F_1,\ldots,F_d)$ będzie polem rzeczywistym,
którego każda składowa spełnia założenia
twierdzenia~\ref{thm:pde7-jump}. Wtedy
\[
 \operatorname{div}T_F
 =T_{(\operatorname{div}F)_{\mathrm{reg}}}
 +[F]\cdot\nu\,\delta_S,\qquad [F]=F_+-F_-.
\]
Równanie $\operatorname{div}T_F=0$ jest równoważne
warunkom $\operatorname{div}F=0$ poza $S$ oraz
$[F]\cdot\nu=0$ na $S$.
\end{wniosek}
\begin{proof}
Sumujemy wzór~\eqref{eq:pde7-jump} dla składowych
$F_j$ i pochodnych $\partial_j$.
Jeśli $\operatorname{div}T_F=0$, testy o nośnikach
w $\Omega\setminus S$ dają zerowanie klasycznej
dywergencji prawie wszędzie na tym zbiorze na mocy
stwierdzenia~\ref{prop:pde7-injective}; z ciągłości
jest ona zerowa punktowo. Ponieważ $\lambda_d(S)=0$,
pozostaje $b\delta_S=0$, gdzie $b=[F]\cdot\nu$ jest
ciągłą funkcją rzeczywistą na $S$.
Gdyby $b(s_0)\ne0$, w dostatecznie małym otoczeniu
$s_0$ funkcja $b$ miałaby stały znak i moduł
ograniczony od dołu przez dodatnią stałą.
Wybierając nieujemną funkcję testową o nośniku
w tym otoczeniu, dodatnią w otoczeniu $s_0$,
otrzymalibyśmy niezerową całkę powierzchniową:
każdy niepusty otwarty fragment wykresu ma dodatnią
miarę powierzchniową (dla $d=1$ miarę liczącą).
Sprzeczność dowodzi $b=0$. Odwrotnie, gdy oba
warunki zachodzą, oba składniki wzoru znikają.
Jest to mechanizm stojący za warunkiem
Rankine'a--Hugoniota z rozdziału~\ref{chap:pde-characteristics}.
\end{proof}

\section{Nośnik, lokalność i zbieżność dystrybucji}
\label{sec:pde7-support}

Dla otwartego $U\subset\Omega$ każdą funkcję
$\varphi\in\mathcal D(U)$ przedłużamy przez zero
do $\Omega$. Otrzymujemy funkcję testową w $\Omega$,
ponieważ nośnik jest zwartym podzbiorem $U$.
Ograniczenie dystrybucji $T$ do $U$ definiujemy przez
$\langle T|_U,\varphi\rangle=\langle T,\widetilde\varphi\rangle$,
gdzie $\widetilde\varphi$ jest tym przedłużeniem.
Dla każdego $K\Subset U$ półnormy funkcji
$\varphi\in\mathcal D_K(U)$ i jej przedłużenia na $K$
są równe. Zatem lokalne oszacowanie
~\eqref{eq:pde7-estimate} dla $T$ daje to samo
oszacowanie dla $T|_U$. W szczególności
$T|_U\in\mathcal D'(U)$.

\begin{definicja}[Nośnik dystrybucji]
Dla $T\in\mathcal D'(\Omega)$ nośnik $\supp T$ jest
dopełnieniem w $\Omega$ sumy wszystkich otwartych
zbiorów $U\subset\Omega$ takich, że $T|_U=0$.
Jest więc względnie domknięty w $\Omega$.
\end{definicja}

Uzasadnijmy, że dystrybucja znika również na tej sumie.
Jeżeli $\varphi$ jest funkcją testową na sumie
otwartych zbiorów, na których $T$ znika, to jej
zwarty nośnik ma skończone pokrycie takimi zbiorami.
Rozkładamy $\varphi$ na skończoną sumę funkcji
testowych $\varphi_i$ o nośnikach w odpowiednich
zbiorach pokrycia, stosując konstrukcję z dowodu
twierdzenia~\ref{thm:pde7-jump} (funkcje odcinające
na kulach, a następnie dzielenie przez ich sumę
dodatnią w otoczeniu nośnika).
Każdy składnik spełnia $\langle T,\varphi_i\rangle=0$,
zatem $\langle T,\varphi\rangle=0$.
W konsekwencji $\Omega\setminus\supp T$ jest
największym otwartym zbiorem, na którym $T$ znika.
W szczególności dla testu o nośniku rozłącznym
z $\supp T$ mamy $\langle T,\varphi\rangle=0$.
Ten sam argument pokazuje lokalność równości:
jeżeli $T_1|_{U_i}=T_2|_{U_i}$ na każdym zbiorze
otwartego pokrycia $\Omega$, to $T_1=T_2$.

Dla $f\in L^1_{\mathrm{loc}}(\Omega)$ zachodzi
\[
 \supp T_f=\Omega\setminus
 \bigcup\{U\subset\Omega: U\text{ otwarty},\
 f=0\text{ prawie wszędzie na }U\}.
\]
Istotnie, jeśli $f=0$ prawie wszędzie na $U$,
to $T_f|_U=0$ bezpośrednio z definicji.
Odwrotna implikacja wynika z iniektywności
odwzorowania $f\mapsto T_f$ na $U$
(stwierdzenie~\ref{prop:pde7-injective}).
Prawa strona wzoru to nośnik istotny funkcji $f$;
nie zależy on od reprezentanta klasy prawie wszędzie.
Ponadto na otwartej sumie występującej we wzorze
$f$ znika prawie wszędzie: z drugiej przeliczalności
można wybrać przeliczalne podpokrycie tej sumy.

Jeżeli $T$ znika na $U$, to dla każdego testu na $U$
jego pochodne również są testami na $U$.
Wzór~\eqref{eq:pde7-distributional-derivative} daje
więc $(\partial^\alpha T)|_U=0$. W szczególności
$\supp(\partial^\alpha T)\subset\supp T$.
Podobnie, jeśli $a\in C^\infty(\Omega)$, to
$\supp(aT)\subset\supp a\cap\supp T$:
na dopełnieniu nośnika $T$ funkcjonał znika,
a na dopełnieniu nośnika $a$ sam iloczyn
$a\varphi$ jest zerowy. Stosujemy tu wykazaną
własność znikania na sumie zbiorów otwartych.

Na przykład $\supp\delta_a=\{a\}$, ponieważ
delta znika poza $a$, natomiast w każdym
otoczeniu $a$ istnieje funkcja testowa
o wartości $1$ w $a$. Również
$\supp(\partial^\alpha\delta_a)=\{a\}$:
inkluzja wynika z powyższego, a test postaci
wielomian razy funkcja odcinająca, skonstruowany
w sekcji~\ref{sec:pde7-derivatives}, daje
niezerową wartość funkcjonału w dowolnym
otoczeniu $a$.

\begin{definicja}[Zbieżność dystrybucji]
Mówimy, że ciąg $T_n\in\mathcal D'(\Omega)$ zbiega do
$T\in\mathcal D'(\Omega)$ w sensie dystrybucji, jeśli
\[
 \langle T_n,\varphi\rangle\longrightarrow
 \langle T,\varphi\rangle
 \qquad\text{dla każdej }\varphi\in\mathcal D(\Omega).
\]
Jest to zbieżność słaba* względem pary
$\mathcal D'(\Omega),\mathcal D(\Omega)$.
\end{definicja}

Warto sprawdzić, że granica punktowa funkcjonałów
liniowych w tej definicji rzeczywiście jest
dystrybucją, jeśli wszystkie granice skalarne istnieją.

\begin{stwierdzenie}[Granica punktowa dystrybucji]
Niech $T_n\in\mathcal D'(\Omega)$ i załóżmy, że dla
każdej $\varphi\in\mathcal D(\Omega)$ istnieje skończona
granica $\lim_{n\to\infty}\langle T_n,\varphi\rangle$.
Wówczas wzór
$\langle T,\varphi\rangle=\lim_{n\to\infty}
\langle T_n,\varphi\rangle$ definiuje dystrybucję.
Co więcej, dla każdego $K\Subset\Omega$ istnieją
$m\in\N_0$ i $C>0$ niezależne od $n$ takie, że
\[
 |\langle T_n,\varphi\rangle|\le C p_{K,m}(\varphi)
 \qquad(\varphi\in\mathcal D_K(\Omega),\ n\ge1).
\]
\end{stwierdzenie}
\begin{proof}
Najpierw uzasadnijmy zupełność $\mathcal D_K(\Omega)$
w topologii jej półnorm. Dla niepustego $K$ topologię
tę wyznacza metryka
\[
 d_K(\varphi,\psi)=\sum_{m=0}^\infty 2^{-m-1}
 \min\{1,p_{K,m}(\varphi-\psi)\}.
\]
Nierówność trójkąta wynika z podaddytywności
$\min\{1,t\}$ dla $t\ge0$ i półnorm, a rozdzielanie
punktów z półnormy $p_{K,0}$.
Zbieżność w tej metryce jest równoważna
zbieżności we wszystkich półnormach: w jedną
stronę szacujemy każdy składnik szeregu,
a w drugą najpierw odcinamy jego mały ogon,
następnie szacujemy skończenie wiele składników.
Tak samo warunek Cauchy'ego w metryce jest
równoważny warunkowi Cauchy'ego w każdej półnormie.

Niech $(\psi_\ell)$ będzie takim ciągiem.
Przedłużamy funkcje przez zero do $\R^d$.
Dla każdego wielowskaźnika $\alpha$ ciąg
$\partial^\alpha\psi_\ell$ jest jednostajnie
Cauchy'ego na $\R^d$, ponieważ wszystkie
pochodne znikają poza $K$. Ma więc jednostajną
granicę $g_\alpha$, ciągłą i znikającą poza $K$.
Dla $x\in\R^d$, $t\in\R$ i wektora standardowej
bazy $e_j$ mamy
\[
 \partial^\alpha\psi_\ell(x+te_j)
 -\partial^\alpha\psi_\ell(x)
 =\int_0^t
 \partial^{\alpha+e_j}\psi_\ell(x+se_j)\,
 \mathrm d\lambda_1(s)
\]
(z całką zorientowaną, gdy $t<0$).
Jednostajna zbieżność pozwala przejść do granicy
po obu stronach. Otrzymujemy tę samą tożsamość
dla $g_\alpha$ i $g_{\alpha+e_j}$.
Z ciągłości funkcji podcałkowej wynika
$\partial_jg_\alpha=g_{\alpha+e_j}$.
Indukcyjnie $g_0$ jest gładka, a jej pochodnymi
są $g_\alpha$. Jej nośnik zawiera się w $K$,
więc $g_0|_\Omega\in\mathcal D_K(\Omega)$,
a ciąg $\psi_\ell$ zbiega do tej funkcji
we wszystkich półnormach. Metryka jest zupełna.
Dla pustego $K$ zupełność jest trywialna.

Ustalmy teraz $K\Subset\Omega$.
Z istnienia granic skalarnych wynika punktowa
ograniczoność rodziny funkcjonałów:
dla każdej $\varphi\in\mathcal D_K(\Omega)$
liczby $\langle T_n,\varphi\rangle$ są ograniczone.
Definiujemy domknięte zbiory
\[
 A_j=\{\varphi\in\mathcal D_K(\Omega):
 |\langle T_n,\varphi\rangle|\le j
 \text{ dla wszystkich }n\ge1\},\qquad j\ge1.
\]
Są domknięte dzięki ciągłości każdego $T_n$ na
$\mathcal D_K(\Omega)$, a ich suma pokrywa
całą przestrzeń $\mathcal D_K(\Omega)$.
Z twierdzenia Baire'a dla przestrzeni metrycznej
zupełnej (jawnie wykorzystywanego tutaj wyniku
topologii) pewien $A_j$ ma niepuste wnętrze.
Istnieją więc $\varphi_0\in A_j$, $m\in\N_0$
i $\delta>0$ takie, że
$\varphi_0+\psi\in A_j$ dla $p_{K,m}(\psi)<\delta$.
Z liniowości mamy wówczas
\[
 |\langle T_n,\psi\rangle|
 \le |\langle T_n,\varphi_0+\psi\rangle|
     +|\langle T_n,\varphi_0\rangle|\le2j.
\]
Dla $p_{K,m}(\varphi)>0$ podstawiamy
$\psi=\delta\varphi/(2p_{K,m}(\varphi))$.
Otrzymujemy
$|\langle T_n,\varphi\rangle|
\le(4j/\delta)p_{K,m}(\varphi)$ dla każdego $n$.
Jeśli półnorma jest zerowa, to $\varphi=0$,
więc oszacowanie również zachodzi.
Dla pustego $K$ wybieramy dowolne $C>0$ i $m=0$.

Granica punktowa jest liniowa przez liniowość
każdego $T_n$, a przejście do granicy w powyższym
oszacowaniu daje
$|\langle T,\varphi\rangle|\le C p_{K,m}(\varphi)$.
Ponieważ $K$ było dowolne, twierdzenie
~\ref{thm:pde7-local-order} daje $T\in\mathcal D'(\Omega)$.
\end{proof}

\begin{stwierdzenie}[Różniczkowanie a granica dystrybucyjna]
\label{prop:pde7-derivative-limit}
Niech $\Omega\subset\R^d$ będzie otwarta i
$T_n\to T$ w $\mathcal D'(\Omega)$. Dla każdego
wielowskaźnika $\alpha$ zachodzi
$\partial^\alpha T_n\to\partial^\alpha T$ w $\mathcal D'(\Omega)$.
Dla każdej ustalonej $a\in C^\infty(\Omega)$ zachodzi
również $aT_n\to aT$ w tym sensie.
\end{stwierdzenie}
\begin{proof}
Dla każdego $\varphi\in\mathcal D(\Omega)$ mamy
\[
 \begin{aligned}
 \langle\partial^\alpha T_n,\varphi\rangle
 &=(-1)^{|\alpha|}\langle T_n,\partial^\alpha\varphi\rangle\\
 &\longrightarrow
 (-1)^{|\alpha|}\langle T,\partial^\alpha\varphi\rangle
 =\langle\partial^\alpha T,\varphi\rangle.
 \end{aligned}
\]
Ponieważ $\partial^\alpha\varphi$ jest funkcją testową,
przejście do granicy jest uprawnione.
Analogicznie $a\varphi$ jest funkcją testową, więc
$\langle aT_n,\varphi\rangle=\langle T_n,a\varphi\rangle
\to\langle T,a\varphi\rangle=\langle aT,\varphi\rangle$.
\end{proof}

Jeśli $f_n\to f$ w $L^1_{\mathrm{loc}}(\Omega)$,
to dla każdego testu o nośniku zawartym w $K$
mamy
\[
 |\langle T_{f_n}-T_f,\varphi\rangle|
 \le p_{K,0}(\varphi)
 \int_K|f_n-f|\,\mathrm d\lambda_d\longrightarrow0.
\]
Otrzymujemy zbieżność dystrybucyjną, ale odwrotna
implikacja nie zachodzi, jak pokazuje poniższy przykład.

\begin{przyklad}[Gładkie funkcje zbieżne do delty]
Jeśli $\rho\in\mathcal D(\R^d)$ spełnia $\rho\ge0$ i
$\int_{\R^d}\rho(x)\,\mathrm d\lambda_d(x)=1$, a
$\rho_\varepsilon(x)=\varepsilon^{-d}\rho(x/\varepsilon)$
dla $\varepsilon>0$, to
$T_{\rho_\varepsilon}\to\delta_0$ w $\mathcal D'(\R^d)$
przy $\varepsilon\downarrow0$.
Rzeczywiście, po podstawieniu $x=\varepsilon y$
otrzymujemy
\[
 \begin{aligned}
 |\langle T_{\rho_\varepsilon}-\delta_0,\varphi\rangle|
 &\le \int_{\R^d}\rho(y)
 |\varphi(\varepsilon y)-\varphi(0)|\,\mathrm d\lambda_d(y)\\
 &\le\sup_{y\in\supp\rho}
 |\varphi(\varepsilon y)-\varphi(0)|
 \longrightarrow0.
 \end{aligned}
\]
Ostatnia granica wynika z ciągłości $\varphi$ w zerze
i ograniczoności nośnika $\rho$.
Nie może to być zbieżność do żadnej funkcji
$f\in L^1_{\mathrm{loc}}(\R^d)$ w normach lokalnych $L^1$:
taka zbieżność implikowałaby $T_f=\delta_0$,
co wykluczyliśmy w sekcji~\ref{sec:pde7-def}.
\end{przyklad}

\section{Wygładzanie i przybliżone jedynki}
\label{sec:pde7-mollifiers}

Wybierzmy odtąd $\rho\in C_c^\infty(B_1(0))$,
$\rho\ge0$, o całce równej $1$, i definiujmy
$\rho_\varepsilon(x)=\varepsilon^{-d}\rho(x/\varepsilon)$
dla $\varepsilon>0$. Podstawienie $x=\varepsilon z$
daje $\int_{\R^d}\rho_\varepsilon\,\mathrm d\lambda_d=1$,
a $\supp\rho_\varepsilon\subset B_\varepsilon(0)$.
Dla $u\in L^1_{\mathrm{loc}}(\Omega)$ i punktów
$x\in\Omega_\varepsilon
=\{x\in\Omega:\operatorname{dist}(x,\R^d\setminus\Omega)>\varepsilon\}$
wprowadzamy
$u_\varepsilon(x)=\int_\Omega u(y)\rho_\varepsilon(x-y)
\,\mathrm d\lambda_d(y)$. Dla dystrybucji stosujemy analogiczną
definicję $T_\varepsilon(x)=\langle T(y),\rho_\varepsilon(x-y)\rangle$.
Przyjmujemy odległość od pustego dopełnienia
równą $+\infty$, gdy $\Omega=\R^d$.
Dla $x\in\Omega_\varepsilon$ domknięta kula
$\overline{B_\varepsilon(x)}$ jest zwartym podzbiorem
$\Omega$, więc całka dla $u$ jest bezwzględnie
zbieżna, a argument funkcjonału $T$ jest testem.
W szczególności $(T_u)_\varepsilon=u_\varepsilon$.
Zapis $T(y)$ w nawiasie jest tylko sposobem wskazania zmiennej
funkcji testowej; nie oznacza wartości punktowej dystrybucji.
Splot jest tu \emph{lokalny}: dla danego $x$ musimy mieć
$\operatorname{supp}(\rho_\varepsilon(x-\cdot))\Subset\Omega$.
Nie wolno automatycznie rozciągać takiej definicji na całą
$\R^d$ bez przedłużenia $T$ lub założeń o nośniku.

\begin{twierdzenie}[Wygładzanie dystrybucji]
\label{thm:pde7-mollify}
Niech $\Omega\subset\R^d$ będzie otwarta, $T\in\mathcal D'(\Omega)$
i $\rho_\varepsilon$ jak wyżej. Wówczas
$T_\varepsilon\in C^\infty(\Omega_\varepsilon)$ oraz
\begin{equation}\label{eq:pde7-deriv-mollify}
 \partial^\alpha T_\varepsilon
 =(\partial^\alpha T)*\rho_\varepsilon
 =T*(\partial^\alpha\rho_\varepsilon)
 \quad\text{na }\Omega_\varepsilon.
\end{equation}
Ponadto dla każdej $\varphi\in\mathcal D(\Omega)$
i dostatecznie małego $\varepsilon>0$ zachodzi
\[
 \int_{\Omega_\varepsilon}T_\varepsilon(x)\varphi(x)
 \,\mathrm d\lambda_d(x)
 \longrightarrow\langle T,\varphi\rangle
 \qquad(\varepsilon\downarrow0).
\]
Jest to lokalna zbieżność dystrybucyjna na rosnących
zbiorach $\Omega_\varepsilon$. Gdy $\Omega=\R^d$,
jest to zwykła zbieżność w $\mathcal D'(\R^d)$.
\end{twierdzenie}
\begin{proof}
Ustalmy $\varepsilon>0$ oraz $x_0\in\Omega_\varepsilon$.
Wybierzmy $r>0$ tak, aby
$K=\overline{B_{\varepsilon+r}(x_0)}\Subset\Omega$.
Dla $|x-x_0|<r$ wszystkie testy
$F_x(y)=\rho_\varepsilon(x-y)$ oraz ich pochodne
względem $x$ mają nośniki zawarte w $K$.
Dla dowolnego wielowskaźnika $\alpha$ i wektora
standardowej bazy $e_j$ iloraz różnicowy
\[
 \frac{\partial_x^\alpha F_{x+he_j}
             -\partial_x^\alpha F_x}{h}
\]
zbiega do $\partial_x^{\alpha+e_j}F_x$
we wszystkich półnormach $p_{K,m}$.
Istotnie, po zastosowaniu dowolnej pochodnej
względem $y$ do ilorazu, podstawowe twierdzenie
rachunku całkowego zapisuje go jako średnią
odpowiedniej pochodnej względem $x_j$
na odcinku od $x$ do $x+he_j$.
Jednostajna ciągłość tej pochodnej na rozpatrywanych
zbiorach zwartych daje zbieżność jednostajną w $y$.
Argument bez ilorazu daje również ciągłość
$x\mapsto\partial_x^\alpha F_x$ we wszystkich
półnormach. Lokalna ciągłość $T$ pozwala zatem
indukcyjnie różniczkować działanie funkcjonału:
\[
 \partial_x^\alpha T_\varepsilon(x)
 =\langle T(y),(\partial^\alpha\rho_\varepsilon)(x-y)\rangle.
\]
Wszystkie te pochodne są ciągłe.
Ponieważ
$\partial_y^\alpha\rho_\varepsilon(x-y)
=(-1)^{|\alpha|}(\partial^\alpha\rho_\varepsilon)(x-y)$,
definicja pochodnej dystrybucyjnej daje
\[
 \langle\partial^\alpha T(y),\rho_\varepsilon(x-y)\rangle
 =\langle T(y),(\partial^\alpha\rho_\varepsilon)(x-y)\rangle.
\]
Dwa czynniki $(-1)^{|\alpha|}$ znoszą się.
To dowodzi~\eqref{eq:pde7-deriv-mollify}.

Ustalmy teraz $\varphi\in\mathcal D(\Omega)$,
przedłużoną przez zero na $\R^d$. Dla funkcji
zerowej teza jest oczywista. Dla niezerowej
wybieramy $r>0$ tak, aby
$K=\supp\varphi+\overline{B_{2r}(0)}\Subset\Omega$.
Dla $0<\varepsilon<r$ nośnik $\varphi$
zawiera się w $\Omega_\varepsilon$, a funkcje
$y\mapsto\varphi(x)\rho_\varepsilon(x-y)$
dla wszystkich $x$ mają nośniki w $K$.
Uzasadnijmy zamianę działania $T$ z całką.
Wybieramy domknięty prostokąt w $\R^d$ zawierający
nośnik $\varphi$ w swoim wnętrzu.
Sumy Riemanna całki funkcji
$(x,y)\mapsto\varphi(x)\rho_\varepsilon(x-y)$
względem $x$ zbiegają do całki we wszystkich
półnormach $p_{K,m}$: po każdej pochodnej
względem $y$ błąd jest ograniczony przez objętość
prostokąta razy moduł jednostajnej ciągłości
odpowiedniej pochodnej na iloczynie prostokąta i $K$,
dla średnicy podziału dążącej do zera.
Różniczkowanie samej całki względem $y$
jest dozwolone przez ciągłość i ograniczoność
pochodnych na tym iloczynie.
Zatem granica sum jest funkcją testową z nośnikiem
w $K$. Ciągłość $T$ na $\mathcal D_K(\Omega)$
pozwala przenieść działanie $T$ przez granicę sum.
Po stronie skalarnej są to sumy Riemanna
funkcji $\varphi T_\varepsilon$, przedłużonej przez
zero poza $\Omega_\varepsilon$; jest ona ciągła,
gdyż nośnik $\varphi$ leży wewnątrz tego zbioru.
Otrzymujemy
\[
 \begin{aligned}
 \int_{\Omega_\varepsilon}T_\varepsilon(x)\varphi(x)
 \,\mathrm d\lambda_d(x)
 &=\left\langle T(y),\int_{\R^d}
 \varphi(x)\rho_\varepsilon(x-y)\,\mathrm d\lambda_d(x)
 \right\rangle\\
 &=\langle T,\check\rho_\varepsilon*\varphi\rangle,
 \end{aligned}
\]
gdzie $\check\rho(z)=\rho(-z)$, a
$\check\rho_\varepsilon(z)=\rho_\varepsilon(-z)$.
Dla każdego wielowskaźnika $\alpha$ zachodzi
\[
 \partial^\alpha(\check\rho_\varepsilon*\varphi)(y)
 =\int_{\R^d}\rho(z)\partial^\alpha\varphi(y+\varepsilon z)
 \,\mathrm d\lambda_d(z).
\]
Jednostajna ciągłość $\partial^\alpha\varphi$
na $\R^d$, zwarty nośnik $\rho$ i jej całka $1$
dają zbieżność jednostajną do $\partial^\alpha\varphi$.
Wszystkie nośniki dla $\varepsilon<r$ są zawarte w $K$.
Z charakterystyki zbieżności w $\mathcal D(\Omega)$
wynika $\check\rho_\varepsilon*\varphi\to\varphi$,
a ciągłość $T$ daje tezę.
\end{proof}

\begin{twierdzenie}[Aproksymacja w $L^p_{\mathrm{loc}}$]
\label{thm:pde7-approx-id}
Niech $\Omega\subset\R^d$ będzie otwarta,
$1\le p<\infty$ i $u\in L^p_{\mathrm{loc}}(\Omega)$.
Wówczas $u_\varepsilon\to u$ w $L^p(K)$ dla każdego
$K\Subset\Omega$.
\end{twierdzenie}
\begin{proof}
Wystarczy rozważyć niepusty $K\Subset\Omega$.
Wybieramy $r>0$ tak, aby
$K' = K+\overline{B_{2r}(0)}\Subset\Omega$.
Niech $v=u\chi_{K'}$ na $\Omega$ i $v=0$
poza $\Omega$. Wtedy $v\in L^p(\R^d)$, a dla
$x\in K$, $0<\varepsilon<r$, zachodzi
$u_\varepsilon(x)=(v*\rho_\varepsilon)(x)$,
gdyż w całce występują tylko punkty z $K'$.
Ponadto $u=v$ prawie wszędzie na $K$.
Podstawienie $y=x-\varepsilon z$ i normalizacja
jądra dają różnicę jako całkę przesunięć $v$
minus $v$. Z całkowej nierówności Minkowskiego
otrzymujemy
\[
 \|v*\rho_\varepsilon-v\|_{L^p(\R^d)}
 \le\int_{\R^d}\rho(z)
 \|v(\cdot-\varepsilon z)-v\|_{L^p(\R^d)}
 \,\mathrm d\lambda_d(z).
\]
Uzasadnijmy ciągłość przesunięć bez odwołania
do gęstości funkcji gładkich, która sama może
wynikać z twierdzenia o wygładzaniu.
Funkcje ciągłe o zwartym nośniku są gęste w $L^p$
dla $p<\infty$: funkcje proste o nośniku skończonej
miary są gęste, a regularność miary pozwala
aproksymować ich składniki charakterystyczne
funkcjami ciągłymi o zwartym nośniku.
Dokładniej, dla zbioru $A$ skończonej miary
wybieramy zwarty $K_0\subset A$ i otwarty $O\supset A$
z $\lambda_d(O\setminus K_0)<\gamma$.
Funkcja ciągła o wartościach w $[0,1]$,
równa $1$ na $K_0$ i o zwartym nośniku w $O$,
może być skonstruowana za pomocą odległości
od $K_0$, jak w dowodzie twierdzenia Fejéra
w rozdziale o szeregach Fouriera.
Dla pustego $K_0$ bierzemy funkcję zerową.
Błąd w normie $L^p$ jest co najwyżej $\gamma^{1/p}$.
Nierówność trójkąta daje żądaną gęstość.

Dla $w\in C_c(\R^d)$ różnice $w(\cdot-b)-w$
przy $|b|\le1$ mają wspólny zwarty nośnik,
a ich normy supremum dążą do zera przy $b\to0$.
Stąd również ich normy $L^p$ dążą do zera.
Jeśli $\|v-w\|_p<\gamma$, to izometryczność
przesunięć daje
\[
 \|v(\cdot-b)-v\|_p
 \le 2\gamma+\|w(\cdot-b)-w\|_p.
\]
Z dowolności $\gamma$ wynika ciągłość przesunięć.
Dla $z\in\supp\rho$ mamy $|\varepsilon z|\le\varepsilon$,
więc normy pod całką w oszacowaniu Minkowskiego
dążą do zera jednostajnie względem $z$.
Ich całka dąży do zera, ponieważ $\int\rho=1$.
W konsekwencji
\[
 \|u_\varepsilon-u\|_{L^p(K)}
 \le\|v*\rho_\varepsilon-v\|_{L^p(\R^d)}
 \longrightarrow0.
\]
Dowodzi to tezy.
\end{proof}

\begin{uwaga}[Nie mylić sposobów zbieżności]
Dla $p=\infty$ zbieżność w normie na ogół nie zachodzi.
Na przykład dla funkcji Heaviside'a $H$ na $\R$
i każdej funkcji ciągłej $v$ na $[-1,1]$ mamy
$\|v-H\|_{L^\infty(-1,1)}\ge1/2$.
Istotnie, jeśli ta norma jest równa $M<\infty$,
to przez ciągłość $|v(x)|\le M$ dla wszystkich
$x<0$ w przedziale, a $|v(x)-1|\le M$ dla $x>0$.
Przechodząc do zera z obu stron, otrzymujemy
$1\le |v(0)|+|1-v(0)|\le2M$.
W szczególności nierówność dotyczy każdego
wygładzenia $H_\varepsilon$.
Dla ciągłej funkcji $u$ na $\Omega$ sytuacja
jest inna: na każdym zwartym $K\Subset\Omega$
dla małych $\varepsilon$ mamy
\[
 \sup_{x\in K}|u_\varepsilon(x)-u(x)|
 \le\sup_{\substack{x\in K\\|z|\le1}}
 |u(x-\varepsilon z)-u(x)|\longrightarrow0,
\]
z jednostajnej ciągłości $u$ na zwartym
otoczeniu $K$.
\end{uwaga}

\begin{center}
\begin{tikzpicture}[scale=.86,>=Stealth]
 \draw[->] (-4.3,0)--(4.7,0) node[below] {$x$};
 \draw[->] (0,-.1)--(0,3.5) node[left] {$\rho_\varepsilon(x)$};
 \draw[thick,blue!65!black,domain=-1.998:1.998,samples=101,smooth]
   plot (\x,{1.35*exp(1-1/(1-(\x/2)^2))});
 \draw[thick,red!70!black,domain=-.999:.999,samples=101,smooth]
   plot (\x,{2.7*exp(1-1/(1-(\x)^2))});
 \draw[blue!65!black,densely dashed] (-2,0)--(-2,-.2) (2,0)--(2,-.2);
 \draw[red!70!black,densely dashed] (-1,0)--(-1,-.2) (1,0)--(1,-.2);
 \node[below] at (-2,-.2) {$-2$};
 \node[below] at (-1,-.2) {$-1$};
 \node[below] at (1,-.2) {$1$};
 \node[below] at (2,-.2) {$2$};
 \node[blue!65!black,anchor=west] at (2.05,1.05) {$\varepsilon=2$};
 \node[red!70!black,anchor=west] at (.75,2.65) {$\varepsilon=1$};
 \node[below] at (0,-.1) {$0$};
\end{tikzpicture}
\par\smallskip\small
Jednowymiarowe profile znormalizowanego jądra
$\rho(x)=c\eta(x)$ z sekcji~\ref{sec:pde7-tests},
gdzie $c=(\int_{\R}\eta\,\mathrm d\lambda_1)^{-1}$.
Wspólna umowna skala pionowa: profil dla
$\varepsilon=1$ jest dwukrotnie wyższy i dwukrotnie
węższy od profilu dla $\varepsilon=2$.
Obie funkcje mają całkę $1$ i znikają poza
odpowiednimi przedziałami nośników.
\end{center}

\section{Słabe pochodne i pierwsze przestrzenie Sobolewa}
\label{sec:pde7-weak}

\begin{definicja}[Pochodna słaba]
\label{def:pde7-weak}
Niech $\Omega\subset\R^d$ będzie otwarta,
$u\in L^1_{\mathrm{loc}}(\Omega)$ i $\alpha\in\N_0^d$.
Funkcję $v\in L^1_{\mathrm{loc}}(\Omega)$ nazywamy
\emph{słabą pochodną} $\partial^\alpha u$, jeżeli
\begin{equation}\label{eq:pde7-weak}
 \int_\Omega u(x)\partial^\alpha\varphi(x)\,\mathrm d\lambda_d(x)
 =(-1)^{|\alpha|}\int_\Omega v(x)\varphi(x)\,
 \mathrm d\lambda_d(x)
 \quad (\varphi\in\mathcal D(\Omega)).
\end{equation}
\end{definicja}

Z definicji pochodnej dystrybucyjnej
~\eqref{eq:pde7-distributional-derivative}
wynika, że~\eqref{eq:pde7-weak} jest równoważne
równości $T_v=\partial^\alpha T_u$:
wystarczy pomnożyć obie strony tożsamości całkowej
przez $(-1)^{|\alpha|}$.
Jeśli $v_1,v_2$ spełniają tę tożsamość dla tego samego
$u$ i $\alpha$, to $T_{v_1-v_2}=0$. Ze stwierdzenia
~\ref{prop:pde7-injective} wynika
$v_1=v_2$ prawie wszędzie. Słaba pochodna jest więc
jednoznacznie wyznaczoną klasą funkcji modulo równość
prawie wszędzie. Zmiana reprezentanta $u$ nie wpływa
na tożsamość~\eqref{eq:pde7-weak}.

Jeśli $u\in C^k(\Omega)$, to dla $|\alpha|\le k$
klasyczna pochodna $\partial^\alpha u$ jest słabą
pochodną: stosujemy kolejno stwierdzenie
~\ref{prop:pde7-classical}.
Jeśli natomiast $v$ jest słabą pochodną $\partial^\alpha u$
i $w$ jest słabą pochodną $\partial^\beta v$, to
\[
 T_w=\partial^\beta T_v
 =\partial^\beta\partial^\alpha T_u
 =\partial^{\alpha+\beta}T_u.
\]
Zatem $w$ jest słabą pochodną $\partial^{\alpha+\beta}u$.
Nie wynika stąd istnienie pośrednich słabych pochodnych
z samego istnienia pochodnej wyższego rzędu.

Pochodna dystrybucyjna istnieje zawsze, ale słaba
pochodna w powyższym znaczeniu istnieje tylko wtedy,
gdy odpowiadająca jej dystrybucja jest regularna.
To rozróżnienie jest istotne: pochodna $H$ jest deltą,
która nie pochodzi od funkcji lokalnie całkowalnej.

\begin{definicja}[Przestrzeń $W^{1,p}$]
Niech $\Omega\subset\R^d$ będzie otwarta i $1\le p\le\infty$.
Przez $W^{1,p}(\Omega)$ oznaczamy zbiór tych
$u\in L^p(\Omega)$, dla których wszystkie pochodne
dystrybucyjne $\partial_jT_u$ są reprezentowane przez
funkcje z $L^p(\Omega)$. Przestrzeń wyposażamy w normę
$\|u\|_{W^{1,p}}=\|u\|_{L^p}+\sum_{j=1}^d
\|\partial_ju\|_{L^p}$.
\end{definicja}

Każda funkcja z $L^p(\Omega)$ dla $1\le p\le\infty$
jest lokalnie całkowalna: na zwartym zbiorze $K$
o skończonej mierze stosujemy nierówność Höldera
(dla $p=1$ wystarczy ograniczenie całki, a dla
$p=\infty$ oszacowanie przez $\lambda_d(K)\|u\|_\infty$).
Liniowość pochodnych dystrybucyjnych i ich jednoznaczność
dają liniowość słabych pochodnych na $W^{1,p}(\Omega)$.
Wobec tego jest to przestrzeń wektorowa, a podany
wzór rzeczywiście definiuje normę: jednorodność
i nierówność trójkąta wynikają z odpowiednich
własności norm $L^p$, a zerowanie normy wymusza
$u=0$ prawie wszędzie już z pierwszego składnika.
Pełną teorię przestrzeni Sobolewa, włącznie ze
zupełnością i twierdzeniami o zanurzeniach,
rozwiniemy w rozdziale~\ref{chap:pde-sobolev}.

\begin{uwaga}[Jak ten skrót ma się do dalszych rozdziałów?]
W tym rozdziale przestrzeń $W^{1,p}(\Omega)$ pojawia się
jedynie jako pierwszy przykład funkcji o słabych pochodnych.
Rozdział~\ref{chap:pde-sobolev} zbuduje ogólną teorię
$W^{k,p}(\Omega)$ na obszarach, wraz z zupełnością,
aproksymacją i przedłużaniem. Z kolei
rozdział~\ref{chap:pde-fourier-Rd} wprowadza
przestrzenie $H^s(\R^d)$ definiowane przez transformatę Fouriera.
Dla całkowitych rzędów $m\ge0$ na całej przestrzeni
zachodzi $H^m(\R^d)=W^{m,2}(\R^d)$ jako równość
przestrzeni z równoważnymi normami; uzasadnimy
ją w dalszej teorii. Na zbiorach otwartych należy
najpierw ustalić definicję przestrzeni $H^s$;
nie utożsamiamy automatycznie konstrukcji przez
ograniczenia funkcji z całej przestrzeni z definicją
przez słabe pochodne.
\end{uwaga}

\begin{przyklad}[Funkcja z załamaniem i funkcja ze skokiem]
Na $(-1,1)$ funkcja $u(x)=|x|$ należy do
$W^{1,p}((-1,1))$ dla każdego $1\le p\le\infty$:
jej słabą pochodną jest $\operatorname{sgn}(x)$.
Natomiast funkcja Heaviside'a ograniczona do $(-1,1)$
nie należy do $W^{1,p}((-1,1))$ dla żadnego
$1\le p\le\infty$, ponieważ jej pochodna
dystrybucyjna jest deltą i nie ma reprezentanta
lokalnie całkowalnego (dowody w sekcjach
~\ref{sec:pde7-def} i~\ref{sec:pde7-derivatives}).
Dla modułu $x$ zarówno funkcja, jak i jej pierwsza
słaba pochodna są ograniczone na przedziale
skończonej miary. Druga pochodna dystrybucyjna
jest jednak równa $2\delta_0$, więc druga słaba
pochodna nie istnieje w sensie~\eqref{eq:pde7-weak}.
Obie funkcje są niegładkie klasycznie, lecz tylko
pierwsza ma lokalnie całkowalną pierwszą słabą pochodną.
\end{przyklad}

\begin{twierdzenie}[Pochodna zerowa oznacza stałość w wymiarze jeden]
\label{thm:pde7-constant}
Niech $I\subset\R$ będzie niepustym przedziałem otwartym,
$T\in\mathcal D'(I)$ i $T'=0$. Istnieje dokładnie jedna
stała $c\in\C$ taka, że $T=T_c$ na $I$.
\end{twierdzenie}
\begin{proof}
Wybierzmy $\eta\in\mathcal D(I)$ o
$\int_I\eta(x)\,\mathrm d\lambda_1(x)=1$
(znormalizowaną funkcję testową na kuli
zawartej w $I$) i połóżmy $c=\langle T,\eta\rangle$.
Dla dowolnej $\varphi\in\mathcal D(I)$ zdefiniujmy
$\psi=\varphi-\bigl(\int_I\varphi(x)\,\mathrm d\lambda_1(x)\bigr)\eta$.
Wtedy $\int_I\psi(x)\,\mathrm d\lambda_1(x)=0$.
Przedłużamy $\psi$ zerem na $\R$ i wybieramy
domknięty przedział $[a,b]\Subset I$ zawierający
jej nośnik; jest to możliwe, ponieważ $I$ jest
przedziałem, także gdy jest nieograniczony.
Funkcja $\Psi(x)=\int_{-\infty}^x\psi(t)
\,\mathrm d\lambda_1(t)$ jest zerowa dla $x<a$,
a dla $x>b$ całka jest równa całce z $\psi$
po całej prostej, czyli również zeru.
Zatem $\supp\Psi\subset[a,b]$.
Z podstawowego twierdzenia rachunku całkowego
$\Psi'=\psi$ i gładkość $\psi$ daje gładkość
$\Psi$, więc $\Psi|_I\in\mathcal D(I)$.
Mamy $\Psi'=\psi$, a więc
$\langle T,\psi\rangle=\langle T,\Psi'\rangle
=-\langle T',\Psi\rangle=0$. Otrzymujemy
$\langle T,\varphi\rangle=c\int_I\varphi(x)\,\mathrm d\lambda_1(x)$.
Jednoznaczność $c$ wynika z wyboru testu o niezerowej całce.
\end{proof}

\begin{wniosek}[Zerowy gradient dystrybucyjny]
Niech $\Omega\subset\R^d$ będzie niepustym spójnym
zbiorem otwartym, a $T\in\mathcal D'(\Omega)$
spełnia $\partial_jT=0$ dla $j=1,\ldots,d$.
Wówczas istnieje dokładnie jedna stała $c\in\C$
taka, że $T=T_c$.
W szczególności jeśli $u\in L^1_{\mathrm{loc}}(\Omega)$
ma wszystkie pierwsze słabe pochodne równe zeru,
to $u=c$ prawie wszędzie.
Na niespójnym zbiorze otwartym wniosek zachodzi
osobno na każdej składowej spójnej.
\end{wniosek}
\begin{proof}
Wybierzmy dowolną kulę $B=B_r(a)$ z
$\overline{B_{2r}(a)}\Subset\Omega$.
Dla $0<\varepsilon<r$ wygładzenie $T_\varepsilon$
jest gładkie na $B$ i na mocy
~\eqref{eq:pde7-deriv-mollify} mamy
$\partial_jT_\varepsilon=(\partial_jT)*\rho_\varepsilon=0$.
Z całkowania wzdłuż odcinka łączącego dwa
punkty kuli (pozostaje on w kuli) wynika
$T_\varepsilon(x)=c_\varepsilon$ dla $x\in B$.
Wybierzmy test $\eta\in\mathcal D(B)$
o całce $1$. Z twierdzenia~\ref{thm:pde7-mollify}
otrzymujemy
\[
 c_\varepsilon=\int_B T_\varepsilon(x)\eta(x)\,
 \mathrm d\lambda_d(x)\longrightarrow
 c_B:=\langle T,\eta\rangle.
\]
Dla dowolnego $\varphi\in\mathcal D(B)$ to samo
twierdzenie daje
\[
 \langle T,\varphi\rangle
 =\lim_{\varepsilon\downarrow0}
 c_\varepsilon\int_B\varphi\,\mathrm d\lambda_d
 =c_B\int_B\varphi\,\mathrm d\lambda_d.
\]
Zatem $T|_B=T_{c_B}$.

Jeśli dwie takie kule przecinają się, to na ich
otwartym przecięciu wybieramy test o całce $1$.
Porównanie dwóch ograniczeń $T$ daje równość
odpowiadających im stałych. Kule tego rodzaju
pokrywają $\Omega$, zatem wyznaczają dobrze
określoną lokalnie stałą funkcję $c(x)$ na $\Omega$.
Spójność $\Omega$ wymusza jej stałość: zbiór punktów,
w których przyjmuje ustaloną wartość, oraz jego
dopełnienie są otwarte. Z lokalności równości
dystrybucji (sekcja~\ref{sec:pde7-support})
otrzymujemy $T=T_c$ na $\Omega$.
Jednoznaczność wynika z testu o całce $1$.
Dla $T=T_u$ stosujemy jeszcze iniektywność
odwzorowania $u\mapsto T_u$.
Składowe spójne zbioru otwartego w $\R^d$
są otwarte, ponieważ każdy punkt ma spójną kulę
zawartą w tym zbiorze. Możemy zatem stosować
powyższy dowód do każdej z nich.
\end{proof}

\section{Rozwiązania fundamentalne operatorów różniczkowych}
\label{sec:pde7-fundamental}

\begin{definicja}[Rozwiązanie fundamentalne]
Niech $P(\partial)=\sum_{|\alpha|\le m}a_\alpha\partial^\alpha$
będzie operatorem różniczkowym o stałych współczynnikach
na $\R^d$, gdzie $a_\alpha\in\C$.
Dystrybucję $E\in\mathcal D'(\R^d)$ nazywamy
\emph{rozwiązaniem fundamentalnym} $P$, gdy
$P(\partial)E=\delta_0$ w $\mathcal D'(\R^d)$.
\end{definicja}

Jeśli $f\in C_c^\infty(\R^d)$, możemy zdefiniować
$(E*f)(x)=\langle E(y),f(x-y)\rangle$.
Na mocy poniższego lematu funkcja ta jest gładka
i spełnia $P(\partial)(E*f)=f$.
Stąd rozwiązanie fundamentalne daje narzędzie rozwiązywania
niejednorodnego równania na całej przestrzeni.
Nie twierdzimy na tym etapie, że $E*f$ opisuje wszystkie
rozwiązania ani że bez dodatkowych warunków można w ten sposób
rozwiązywać dowolne zadanie brzegowe.


\begin{lemat}[Różniczkowanie splotu z dystrybucją]
\label{lem:pde7-convolution-derivatives}
Niech $d\ge1$, $T\in\mathcal D'(\R^d)$ i
$\psi\in C_c^\infty(\R^d)$. Funkcja
$u(x)=\langle T(y),\psi(x-y)\rangle$ należy do
$C^\infty(\R^d)$ i dla każdego wielowskaźnika
$\alpha\in\N_0^d$ zachodzi
\[
 \partial^\alpha(T*\psi)
 =T*(\partial^\alpha\psi)=(\partial^\alpha T)*\psi.
\]
W szczególności dla operatora o stałych współczynnikach
$P(\partial)$ i rozwiązania fundamentalnego $E$ mamy
$P(\partial)(E*\psi)=\psi$.
\end{lemat}
\begin{proof}
Niech $A\subset\R^d$ będzie niepustym zbiorem zwartym
i niech $B=\supp\psi$. Dla $x\in A$ nośniki
funkcji $y\mapsto\psi(x-y)$ zawierają się w zwartym
zbiorze $A-B=\{x-z:x\in A,\ z\in B\}$.
Dla różniczkowania w otoczeniu $A$ wystarczy
zastąpić $A$ jego zwartym poszerzeniem.
Stosujemy argument z dowodu twierdzenia
~\ref{thm:pde7-mollify}, w którym dla gładkości
splotu i rachunku pochodnych nie używaliśmy
ani dodatniości jądra, ani warunku całki równej $1$.
Dokładniej, po każdej pochodnej względem $y$
ilorazy różnicowe względem $x_j$ zbiegają
jednostajnie na wspólnym zwartym nośniku, dzięki
podstawowemu twierdzeniu rachunku całkowego
i jednostajnej ciągłości pochodnych $\psi$.
Zatem zbieżność zachodzi we wszystkich półnormach
przestrzeni testowej o tym nośniku.
Ciągłość $T$ i indukcja względem rzędu pochodnej dają
$\partial_x^\alpha u(x)=\langle T(y),
(\partial^\alpha\psi)(x-y)\rangle$,
a ciągłość tych pochodnych wynika z ciągłości
przesunięć testów w tych samych półnormach.
Przypadek $\psi=0$ jest oczywisty.
Dla pochodnej pierwszego rzędu korzystamy z definicji
pochodnej dystrybucyjnej:
\[
 \langle\partial_{y_j}T(y),\psi(x-y)\rangle
 =-\langle T(y),\partial_{y_j}\psi(x-y)\rangle
 =\langle T(y),\partial_{x_j}\psi(x-y)\rangle.
\]
Iteracja daje wszystkie wielowskaźniki. Na koniec
$P(\partial)(E*\psi)=(P(\partial)E)*\psi
=\delta_0*\psi=\psi$, gdzie ostatnia równość
wynika bezpośrednio z działania delty na test.
\end{proof}

\begin{twierdzenie}[Rozwiązanie fundamentalne operatora $-\Delta$]
\label{thm:pde7-laplace}
Niech $d\ge1$ i niech $\omega_d$ oznacza pole powierzchni
sfery jednostkowej $S^{d-1}\subset\R^d$ dla $d\ge2$.
Funkcja lokalnie całkowalna
\begin{equation}\label{eq:pde7-laplace-fundamental}
 \Phi_d(x)=\begin{cases}
 -|x|/2,&d=1,\\
 (2\pi)^{-1}\log(1/|x|),&d=2,\\
 ((d-2)\omega_d)^{-1}|x|^{2-d},&d\ge3,
 \end{cases}
\end{equation}
spełnia $-\Delta T_{\Phi_d}=\delta_0$ na $\R^d$.
\end{twierdzenie}
\begin{proof}
Wartość $\Phi_d(0)$ w wymiarach $d\ge2$ ustalamy
dowolnie, na przykład równą $0$.
W wymiarze jeden $\Phi_1$ jest lokalnie całkowalna
i stwierdzenie wynika z $\partial_x^2T_{|x|}=2\delta_0$
(przykład~\ref{ex:pde7-heaviside}).

Dla $d\ge3$ całkowanie w zmiennych biegunowych daje
dla $r>0$
\[
 \int_{B_r(0)}|\Phi_d(x)|\,\mathrm d\lambda_d(x)
 =\frac1{d-2}\int_0^r s\,\mathrm d\lambda_1(s)
 =\frac{r^2}{2(d-2)}.
\]
Dla $d=2$ mamy $\omega_2=2\pi$ i dla $0<r<1$
\[
 \begin{aligned}
 \int_{B_r(0)}|\Phi_2(x)|\,\mathrm d\lambda_2(x)
 &=\int_0^r s\log(1/s)\,\mathrm d\lambda_1(s)\\
 &=\frac{r^2}{2}\log(1/r)+\frac{r^2}{4}.
 \end{aligned}
\]
Granica dolnego wyrazu pierwotnej jest zerowa,
gdyż $s^2\log(1/s)\to0$ dla $s\downarrow0$.
W obu przypadkach całka dąży do zera przy $r\downarrow0$.
Z gładkości poza zerem wynika więc lokalna
całkowalność na całej $\R^d$.

Niech $F(r)$ oznacza profil radialny $\Phi_d$ dla $r>0$.
Bezpośrednie różniczkowanie obu wzorów daje
\[
 F'(r)=-\frac1{\omega_dr^{d-1}},\qquad
 F''(r)=\frac{d-1}{\omega_dr^d}.
\]
Dla $x\ne0$, $r=|x|$,
$\partial_jr=x_j/r$ oraz
$\partial_j^2r=1/r-x_j^2/r^3$.
Reguła łańcuchowa daje więc
\[
 \Delta\Phi_d(x)
 =F''(r)+\frac{d-1}{r}F'(r)=0.
\]
Dla $\varphi\in\mathcal D(\R^d)$ wybieramy $R>0$
tak, aby $\supp\varphi\Subset B_R(0)$.
Dla $0<\varepsilon<R$ stosujemy drugą tożsamość
Greena~\eqref{eq:pde3-green2} na pierścieniu
$A_{\varepsilon,R}=B_R(0)\setminus\overline{B_\varepsilon(0)}$
(do części rzeczywistej i urojonej testu osobno).
Ponieważ $\Delta\Phi_d=0$ na pierścieniu,
otrzymujemy
\[
 -\int_{A_{\varepsilon,R}}\Phi_d\Delta\varphi
 \,\mathrm d\lambda_d
 =\int_{\partial A_{\varepsilon,R}}
 (\varphi\partial_\nu\Phi_d-\Phi_d\partial_\nu\varphi)
 \,\mathrm d\sigma.
\]
Na zewnętrznej sferze znikają test i jego pochodne.
Na wewnętrznej sferze $\nu=-x/|x|$, zatem
$\partial_\nu\Phi_d=-F'(\varepsilon)
=1/(\omega_d\varepsilon^{d-1})$.
W konsekwencji
\[
 -\int_{B_R\setminus B_\varepsilon}\Phi_d(x)\Delta\varphi(x)
 \,\mathrm d\lambda_d(x)
 =\int_{\partial B_\varepsilon}\varphi(x)
 \frac{1}{\omega_d\varepsilon^{d-1}}\,\mathrm d\sigma(x)
 -\int_{\partial B_\varepsilon}\Phi_d(x)
 \partial_\nu\varphi(x)\,\mathrm d\sigma(x).
\]
Miara powierzchniowa sfery o promieniu $\varepsilon$
wynosi $\omega_d\varepsilon^{d-1}$. Po odjęciu
$\varphi(0)$ od pierwszej całki moduł różnicy
jest więc co najwyżej
$\sup_{|x|=\varepsilon}|\varphi(x)-\varphi(0)|$,
co dąży do zera z ciągłości testu.
Moduł drugiej całki nie przekracza
\[
 \|\nabla\varphi\|_\infty\,
 |F(\varepsilon)|\,\omega_d\varepsilon^{d-1}.
\]
Dla $d\ge3$ iloczyn bez normy gradientu wynosi
$\varepsilon/(d-2)$, dla $d=2$ i $\varepsilon<1$
wynosi $\varepsilon\log(1/\varepsilon)$.
Obie wielkości dążą do zera.
Ponadto moduł pominiętej całki po $B_\varepsilon(0)$
jest ograniczony przez
$\|\Delta\varphi\|_\infty
\int_{B_\varepsilon(0)}|\Phi_d|\,\mathrm d\lambda_d$,
co dąży do zera z wcześniejszych obliczeń.
Ze znaku $(-1)^2=1$ przy drugiej pochodnej
dystrybucyjnej otrzymujemy
\[
 \langle-\Delta T_{\Phi_d},\varphi\rangle
 =-\int_{\R^d}\Phi_d(x)\Delta\varphi(x)
 \,\mathrm d\lambda_d(x)=\varphi(0).
\]
Jest to żądana równość.
\end{proof}

\begin{uwaga}[Zgodność z wcześniejszą teorią potencjału]
W rozdziale~\ref{chap:pde-laplace},
w twierdzeniu~\ref{thm:pde3-phi}, otrzymaliśmy te same
funkcje dla $d\ge2$ oraz tę samą tożsamość całkową.
Tutaj interpretujemy ją w przestrzeni dystrybucji
i uwzględniamy także przypadek $d=1$.
Z lematu~\ref{lem:pde7-convolution-derivatives}
wynika w szczególności, że dla $f\in\mathcal D(\R^d)$
funkcja
\[
 u(x)=\int_{\R^d}\Phi_d(x-y)f(y)\,\mathrm d\lambda_d(y)
\]
jest gładka i spełnia $-\Delta u=f$.
Całka jest bezwzględnie zbieżna dzięki zwartemu
nośnikowi $f$ i lokalnej całkowalności $\Phi_d$;
podstawienie $z=x-y$ utożsamia ją ze splotem
$T_{\Phi_d}*f$. Nie wyciągamy stąd wniosku
o jednoznaczności rozwiązania bez dodatkowych
warunków: dodanie dowolnej gładkiej funkcji
harmonicznej daje kolejne rozwiązanie.
\end{uwaga}

\section{Równanie ciepła i delta w chwili początkowej}
\label{sec:pde7-heat}

\begin{twierdzenie}[Jądro cieplne jako rozwiązanie fundamentalne]
\label{thm:pde7-heat}
Niech $d\ge1$ i $\kappa>0$. Na $\R^{1+d}$ definiujemy
funkcję
\[
 E(t,x)=\begin{cases}
 (4\pi\kappa t)^{-d/2}
 \exp\!\bigl(-|x|^2/(4\kappa t)\bigr),&t>0,\\
 0,&t\le0.
 \end{cases}
\]
Wtedy $E\in L^1_{\mathrm{loc}}(\R^{1+d})$ i dla
operatora $L=\partial_t-\kappa\Delta_x$ zachodzi
$L T_E=\delta_{(0,0)}$ w $\mathcal D'(\R^{1+d})$.
\end{twierdzenie}
\begin{proof}
Oznaczmy przez $\Gamma_\kappa(t,x)$ jądro cieplne
dla $t>0$. Z twierdzenia~\ref{thm:pde4-gauss-basic}
wiemy, że jest ono dodatnie, gładkie, spełnia
$(\partial_t-\kappa\Delta_x)\Gamma_\kappa=0$ oraz
ma całkę przestrzenną równą $1$ dla każdego $t>0$.
Dla dowolnego zwartego zbioru
$K\subset\R^{1+d}$ wybieramy $A>0$ tak, aby
$K\subset[-A,A]\times\R^d$.
Z twierdzenia Tonellego i normalizacji jądra
otrzymujemy
\[
 \int_K|E(t,x)|\,\mathrm d\lambda_{1+d}(t,x)
 \le\int_0^A\int_{\R^d}\Gamma_\kappa(t,x)
 \,\mathrm d\lambda_d(x)\,\mathrm d\lambda_1(t)=A.
\]
Zatem $E$ jest lokalnie całkowalna mimo osobliwości
w punkcie $(0,0)$.

Ustalmy $\varphi\in\mathcal D(\R^{1+d})$.
Z definicji pochodnej dystrybucyjnej (znak minus
dla pierwszej pochodnej i znak plus dla drugiej)
wynika
\[
 \langle LT_E,\varphi\rangle
 =\int_0^\infty\int_{\R^d}\Gamma_\kappa(t,x)
 [-\partial_t\varphi(t,x)-\kappa\Delta_x\varphi(t,x)]
 \,\mathrm d\lambda_d(x)\,\mathrm d\lambda_1(t).
\]
Całka jest bezwzględnie zbieżna: pochodne testu
są ograniczone i mają wspólny zwarty nośnik,
a $E$ jest lokalnie całkowalna. Co więcej,
moduł części całki dla $0<t<\varepsilon$ nie
przekracza
\[
 \varepsilon
 \bigl(\|\partial_t\varphi\|_\infty+
             \kappa\|\Delta_x\varphi\|_\infty\bigr),
\]
ponieważ jądro ma masę przestrzenną $1$.
Wystarczy zatem obliczyć całkę $I_\varepsilon$
po $t\ge\varepsilon$ i przejść do granicy.

Wybieramy $A>\varepsilon$ oraz $R>0$ tak, aby
nośnik testu był zawarty w
$(-A,A)\times B_R(0)$. Wszystkie jego pochodne
znikają w otoczeniu górnej granicy czasowej $A$
i przestrzennego brzegu kuli.
Na $[\varepsilon,A]\times\overline{B_R(0)}$
jądro oraz jego pochodne są gładkie i ograniczone.
Możemy więc stosować twierdzenie Fubiniego,
całkowanie przez części w czasie oraz dwukrotne
całkowanie przez części względem przestrzeni.
Otrzymujemy
\[
\begin{aligned}
 I_\varepsilon
 &=\int_{\R^d}\Gamma_\kappa(\varepsilon,x)
           \varphi(\varepsilon,x)\,\mathrm d\lambda_d(x)\\
 &\quad+\int_\varepsilon^A\int_{\R^d}
  (\partial_t\Gamma_\kappa-\kappa\Delta_x\Gamma_\kappa)
       \varphi\,\mathrm d\lambda_d\,\mathrm d\lambda_1.
\end{aligned}
\]
W pierwszym składniku znak jest dodatni,
ponieważ dolna granica całkowania w czasie
występuje z przeciwnym znakiem w całce
z $-\Gamma_\kappa\partial_t\varphi$.
Nie ma składników brzegowych w przestrzeni,
gdyż test wraz z pochodnymi znika przy
brzegu wybranej kuli. Drugi składnik jest
zerowy na mocy równania cieplnego dla $t>0$.

Podstawienie $x=\sqrt{4\kappa\varepsilon}\,z$
o jakobianie $(4\kappa\varepsilon)^{d/2}$ daje
\[
 I_\varepsilon
 =\pi^{-d/2}\int_{\R^d}e^{-|z|^2}
 \varphi(\varepsilon,\sqrt{4\kappa\varepsilon}\,z)
 \,\mathrm d\lambda_d(z).
\]
Dla każdego ustalonego $z$ argument testu zmierza
do $(0,0)$, a moduł funkcji podcałkowej wraz
z czynnikiem normalizacyjnym nie przekracza
$\pi^{-d/2}e^{-|z|^2}\|\varphi\|_\infty$.
Jest to całkowalna majoranta niezależna od
$\varepsilon$. Z twierdzenia o zbieżności
zmajoryzowanej oraz normalizacji całki Gaussa
wynika $I_\varepsilon\to\varphi(0,0)$.
Ostatecznie
\[
 \langle LT_E,\varphi\rangle
 =\lim_{\varepsilon\downarrow0} I_\varepsilon
 =\varphi(0,0)
 =\langle\delta_{(0,0)},\varphi\rangle.
\]
\end{proof}

\begin{uwaga}[Nie jest to rozwiązanie klasyczne w $t=0$]
Funkcja $E$ jest klasycznie gładka dla $t>0$ oraz
dla $t<0$. Jej przedłużenie przez zero na czasy
niedodatnie nie spełnia jednak równania jednorodnego
na całej czasoprzestrzeni: w sensie dystrybucji
punkt $(0,0)$ jest źródłem o masie $1$.
Nie jest nim cała hiperpowierzchnia $t=0$.
W szczególności $LT_E=0$ na każdym otwartym
zbiorze niezawierającym $(0,0)$.

Dla ustalonego $t>0$ jądro przestrzenne
$\Gamma_\kappa(t,\cdot)$ wyznacza dystrybucję na
$\R^d$, która przy $t\downarrow0$ zbiega do
$\delta_0$: dla testu $\psi\in\mathcal D(\R^d)$
stosujemy to samo podstawienie i twierdzenie
o zbieżności zmajoryzowanej do całki
$\int_{\R^d}\Gamma_\kappa(t,x)\psi(x)\,\mathrm d\lambda_d(x)$.
Jest to granica w przestrzeni dystrybucji
przestrzennych, a nie punktowa wartość funkcji
$E(0,x)$.
\end{uwaga}

\begin{wniosek}[Czasoprzestrzenny wzór Duhamela]
Jeśli $f\in\mathcal D(\R^{1+d})$, to funkcja
$u=T_E*f$ jest gładka na całej czasoprzestrzeni,
spełnia $Lu=f$ i ma postać
\[
 u(t,x)=\int_{-\infty}^t\int_{\R^d}
 \Gamma_\kappa(t-s,x-y)f(s,y)
 \,\mathrm d\lambda_d(y)\,\mathrm d\lambda_1(s).
\]
Jeżeli $f(s,y)=0$ dla wszystkich $s<t_0$ i $y\in\R^d$,
to $u(t,x)=0$ dla $t<t_0$.
\end{wniosek}
\begin{proof}
Z lematu~\ref{lem:pde7-convolution-derivatives},
zastosowanego w wymiarze $1+d$, i twierdzenia
~\ref{thm:pde7-heat} wynika
$u\in C^\infty(\R^{1+d})$ oraz
$Lu=(LT_E)*f=\delta_{(0,0)}*f=f$.
Ponieważ $E$ jest lokalnie całkowalna, splot
z testem jest całką bezwzględnie zbieżną:
dla ustalonego $(t,x)$ argument $E$ przebiega
zwarty zbiór będący przesunięciem odbitego
nośnika $f$. Zmiana zmiennych w splocie daje
całkę z $E(t-s,x-y)f(s,y)$.
Dzięki bezwzględnej całkowalności stosujemy
twierdzenie Fubiniego, a ponieważ $E(t-s,x-y)=0$
dla $s\ge t$, pozostaje podany wzór.
Dla $t<t_0$ każda wartość $s<t$ leży w zakresie
zanikania $f$, więc $u(t,x)=0$.
Tak otrzymujemy przyczynowe rozwiązanie równania
z wymuszeniem o zwartym nośniku; nie orzekamy
jednoznaczności bez dalszych założeń.
\end{proof}

\section{Lemat Weyla: rozwiązanie dystrybucyjne jest gładkie}
\label{sec:pde7-weyl}

\begin{twierdzenie}[Lemat Weyla]
\label{thm:pde7-weyl}
Niech $d\ge2$, $\Omega\subset\R^d$ będzie otwarta,
a $T\in\mathcal D'(\Omega)$ spełnia $\Delta T=0$.
Wówczas istnieje dokładnie jedna funkcja
$h\in C^\infty(\Omega;\C)$ taka, że $T=T_h$
oraz $\Delta h=0$ punktowo.
W szczególności $h$ jest rzeczywiście analityczna
(tzn. jej część rzeczywista i urojona są
rzeczywiście analityczne).
\end{twierdzenie}
\begin{proof}
Ustalmy $a\in\Omega$ i wybierzmy $r>0$ takie, że
$\overline{B_{3r}(a)}\Subset\Omega$.
Dla $0<\varepsilon<r$ definiujemy
$u_\varepsilon=T*\rho_\varepsilon$ na $B_{2r}(a)$.
Rzeczywiście, jeśli $x\in B_{2r}(a)$, to
$\overline{B_\varepsilon(x)}\subset B_{3r}(a)$;
w szczególności $x\in\Omega_\varepsilon$.
Z twierdzenia~\ref{thm:pde7-mollify} wynika
$u_\varepsilon\in C^\infty(B_{2r}(a))$, a wzór
~\eqref{eq:pde7-deriv-mollify} daje
\[
 \Delta u_\varepsilon=(\Delta T)*\rho_\varepsilon=0
 \qquad\text{na }B_{2r}(a).
\]

Wybierzmy radialną funkcję
$\eta\in C_c^\infty(B_r(0))$, nieujemną i o całce $1$.
Możemy ją otrzymać przez skalowanie i normalizację
radialnej funkcji testowej z sekcji~\ref{sec:pde7-tests}.
Zapiszmy $\eta(z)=q(|z|)$, gdzie $q$ jest jej
profilem radialnym. Dla $x\in B_r(a)$ i $0<s<r$
mamy $\overline{B_s(x)}\subset B_{2r}(a)$.
Stosując twierdzenie~\ref{thm:pde3-mean}
do części rzeczywistej i urojonej $u_\varepsilon$,
otrzymujemy
\[
 \int_{S^{d-1}}u_\varepsilon(x-s\theta)\,
 \mathrm d\sigma(\theta)
 =\omega_d u_\varepsilon(x).
\]
Wykorzystaliśmy również niezmienniczość miary
powierzchniowej przy przekształceniu
$\theta\mapsto-\theta$.
Po całkowaniu w zmiennych biegunowych dostajemy
\[
\begin{aligned}
 \int_{B_r(0)}u_\varepsilon(x-z)\eta(z)\,
 \mathrm d\lambda_d(z)
 &=\int_0^r q(s)s^{d-1}
 \int_{S^{d-1}}u_\varepsilon(x-s\theta)\,
 \mathrm d\sigma(\theta)\,\mathrm d\lambda_1(s)\\
 &=u_\varepsilon(x)\omega_d
 \int_0^r q(s)s^{d-1}\,\mathrm d\lambda_1(s)
 =u_\varepsilon(x).
\end{aligned}
\]
Ostatnia równość wynika z normalizacji $\eta$.
Wszystkie argumenty $x-z$ występujące w całce
należą do $B_{2r}(a)$, więc nie przedłużamy
$u_\varepsilon$ na całą przestrzeń.
Po podstawieniu $y=x-z$ otrzymujemy na $B_r(a)$
\[
 u_\varepsilon(x)=
 \int_{B_{2r}(a)}u_\varepsilon(y)\eta(x-y)\,
 \mathrm d\lambda_d(y).
\]

Niech teraz $\varphi\in\mathcal D(B_r(a))$;
przedłużamy ją przez zero. Oznaczmy
$\check\eta(z)=\eta(-z)$ (tu $\check\eta=\eta$).
Splot $\check\eta*\varphi$ jest gładki
i ma nośnik zawarty w zwartym zbiorze
$\supp\varphi+\supp\check\eta\subset B_{2r}(a)$.
Gładkość wynika z różniczkowania całki o zwartym
nośniku, a inkluzja nośników z faktu, że poza
wskazaną sumą zbiory występujące w splocie
nie mają punktu wspólnego.
Jest to zatem test w $B_{2r}(a)$.
W całce podwójnej po $x,y$ niezerowe składniki
występują tylko wtedy, gdy $x\in\supp\varphi$
oraz $y\in x-\supp\eta$. Obie zmienne przebiegają
wówczas zbiory zwarte wewnątrz dziedzin
odpowiednich funkcji. Ponieważ $u_\varepsilon$
jest tam ograniczona, całka jest bezwzględnie
zbieżna i można zastosować twierdzenie Fubiniego:
\[
 \int_{B_r(a)}u_\varepsilon(x)\varphi(x)\,
 \mathrm d\lambda_d(x)
 =\int_{B_{2r}(a)}u_\varepsilon(y)
 (\check\eta*\varphi)(y)\,\mathrm d\lambda_d(y).
\]
W twierdzeniu~\ref{thm:pde7-mollify} wykazaliśmy
lokalną zbieżność dystrybucyjną $u_\varepsilon$ do $T$.
Oba testy są ustalone i mają zwarte nośniki
w $\Omega$, więc po przejściu do granicy otrzymujemy
\[
 \langle T,\varphi\rangle
 =\langle T,\check\eta*\varphi\rangle.
\]

Na $B_r(a)$ określamy funkcję
$h_a(x)=\langle T(y),\eta(x-y)\rangle$.
Z argumentu z dowodu twierdzenia
~\ref{thm:pde7-mollify}, z jądrem $\eta$ zamiast
$\rho_\varepsilon$, wynika, że $h_a$ jest gładka.
Argument ten dotyczy także ustalonego jądra:
w otoczeniu każdego $x\in B_r(a)$ wszystkie
przesunięte testy mają wspólny zwarty nośnik
w $B_{2r}(a)$.
Ponadto uzasadniona tam zamiana działania
funkcjonału z całką (przez zbieżność sum Riemanna
we wszystkich półnormach) daje
\[
 \int_{B_r(a)}h_a(x)\varphi(x)\,\mathrm d\lambda_d(x)
 =\langle T,\check\eta*\varphi\rangle
 =\langle T,\varphi\rangle.
\]
Mamy więc $T|_{B_r(a)}=T_{h_a}$.
Różniczkowanie splotu daje również
$\Delta h_a=(\Delta T)*\eta=0$ na tej kuli.

Kule tego rodzaju pokrywają $\Omega$. Jeśli dwie
kule przecinają się, to ich lokalne reprezentacje
wyznaczają tę samą dystrybucję na otwartym
przecięciu. Ze stwierdzenia~\ref{prop:pde7-injective}
są tam równe prawie wszędzie, a dzięki ciągłości
są równe w każdym punkcie (niezerowa różnica
miałaby niezerowy moduł na pewnej kuli
dodatniej miary).
Otrzymujemy zatem dobrze określoną gładką
funkcję $h$ na $\Omega$, która lokalnie jest
równa odpowiednim $h_a$.
Lokalność równości dystrybucji z sekcji
~\ref{sec:pde7-support} daje $T=T_h$,
a lokalna harmoniczność daje $\Delta h=0$.
Jednoznaczność wśród gładkich reprezentacji
wynika ponownie z iniektywności i ciągłości.
Na koniec twierdzenie~\ref{thm:pde3-analytic},
zastosowane osobno do części rzeczywistej
i urojonej, daje analityczność.
\end{proof}

\begin{uwaga}[Znaczenie założenia harmoniczności]
Lemat Weyla nie mówi, że każda dystrybucja jest
funkcją gładką: kluczowym założeniem jest
równanie $\Delta T=0$.
Na przykład rozwiązanie fundamentalne $\Phi_d$
z twierdzenia~\ref{thm:pde7-laplace} jest gładkie
i harmoniczne poza zerem, ale spełnia
$-\Delta T_{\Phi_d}=\delta_0$ na całej przestrzeni.
Nie jest więc rozwiązaniem jednorodnego równania
dystrybucyjnego na żadnym otoczeniu zera.
\end{uwaga}

\section{Ograniczenia rachunku dystrybucji}
\label{sec:pde7-limits}

Pochodne i mnożenie przez funkcje gładkie są dobrze określone
na $\mathcal D'(\Omega)$, jak wykazano w sekcji
~\ref{sec:pde7-derivatives}. Nie można jednak wprowadzić
na przestrzeni wszystkich dystrybucji iloczynu, który
jednocześnie miałby wszystkie zwykłe własności
algebraiczne, zachowywał mnożenie funkcji lokalnie
ograniczonych i spełniał regułę Leibniza.

\begin{stwierdzenie}[Przeszkoda w rozszerzeniu iloczynu]
Na $\mathcal D'(\R)$ nie istnieje przemienny, łączny
i dwuliniowy iloczyn $\star$ spełniający jednocześnie
następujące warunki:
\begin{enumerate}
\item $T_f\star T_g=T_{fg}$ dla wszystkich
$f,g\in L^\infty_{\mathrm{loc}}(\R)$;
\item $(S\star T)'=S'\star T+S\star T'$ dla wszystkich
$S,T\in\mathcal D'(\R)$.
\end{enumerate}
\end{stwierdzenie}
\begin{proof}
Niech $A=T_H$, gdzie $H$ jest funkcją Heaviside'a.
Ponieważ $H^2=H$ prawie wszędzie, pierwszy warunek
daje $A\star A=A$. Z przykładu~\ref{ex:pde7-heaviside}
mamy $A'=\delta_0$. Różniczkując równość
$A\star A=A$ za pomocą drugiego warunku i korzystając
z przemienności, otrzymujemy
\[
 2(A\star\delta_0)=\delta_0.
\]
Mnożymy tę równość przez $A$. Dwuliniowość i łączność
dają
\[
 A\star\delta_0
 =2A\star(A\star\delta_0)
 =2(A\star A)\star\delta_0
 =2A\star\delta_0.
\]
Stąd $A\star\delta_0=0$, a pierwsza równość
wymusza $\delta_0=0$. Jest to sprzeczne
z działaniem delty na test o wartości $1$ w zerze.
\end{proof}

Wynik ten jest konkretną wersją przeszkody Schwartza:
dotyczy dokładnie wymienionych założeń i nie wyklucza
wszystkich częściowo określonych iloczynów ani
innych teorii funkcji uogólnionych.
W szczególności zapis $\delta_0^2$ nie ma
standardowego znaczenia w rachunku dystrybucji.
Również wyrażenie $u\,\partial_xu$ dla nieciągłego
$u$ nie musi mieć sensu: dla $u=H$ wymagałoby
mnożenia delty przez funkcję nieciągłą w zerze.

Dla równań zachowania można ominąć tę trudność,
stosując \emph{postać dywergencyjną}. Dla funkcji
gładkiej reguła łańcuchowa daje
$u\partial_xu=\partial_x(u^2/2)$.
Jeśli rzeczywista funkcja $u$ należy do
$L^2_{\mathrm{loc}}(Q)$, gdzie $Q$ jest otwartym
podzbiorem czasoprzestrzeni $\R^2$, to
$u\in L^1_{\mathrm{loc}}(Q)$ z nierówności
Cauchy'ego--Schwarza na zbiorach zwartych,
a $u^2/2\in L^1_{\mathrm{loc}}(Q)$ z definicji.
Możemy więc nadać precyzyjny sens równaniu
\[
 \partial_tT_u+\partial_xT_{u^2/2}=0
 \quad\text{w }\mathcal D'(Q).
\]
Z definicji pochodnych dystrybucyjnych jest ono
równoważne warunkowi
\[
 \int_Q
 \left(u(t,x)\partial_t\varphi(t,x)
       +\frac{u(t,x)^2}{2}\partial_x\varphi(t,x)\right)
 \,\mathrm d\lambda_2(t,x)=0
 \qquad(\varphi\in\mathcal D(Q)).
\]
Nie wymaga on definiowania iloczynu
$u\,\partial_xT_u$. Sam warunek dystrybucyjny
nie zapewnia jednak jednoznaczności: warunki
entropijne omawiamy w teorii praw zachowania
(por. rozdział~\ref{chap:pde-characteristics}).

\begin{przyklad}[Skąd bierze się problem z $\delta^2$?]
Niech $\rho\in\mathcal D(\R)$ będzie nieujemna
i niech $\int_\R\rho\,\mathrm d\lambda_1=1$.
Dla $\varepsilon>0$ określamy
$\rho_\varepsilon(x)=\varepsilon^{-1}\rho(x/\varepsilon)$.
Wiemy z sekcji~\ref{sec:pde7-support}, że
$T_{\rho_\varepsilon}\to\delta_0$.
Wybieramy rzeczywisty, nieujemny test $\varphi$
taki, że $\varphi(0)>0$.
Podstawienie $x=\varepsilon y$ daje
\[
 \varepsilon\langle T_{\rho_\varepsilon^2},\varphi\rangle
 =\int_\R \rho(y)^2\varphi(\varepsilon y)\,
 \mathrm d\lambda_1(y).
\]
Ponieważ nośnik $\rho$ jest zwarty,
$\varphi(\varepsilon y)\to\varphi(0)$ jednostajnie
na nim. Stąd prawa strona dąży do
\[
 c=\varphi(0)\int_\R\rho(y)^2\,\mathrm d\lambda_1(y)>0.
\]
Dodatniość wynika z $\varphi(0)>0$ oraz z faktu,
że ciągła funkcja $\rho$ o całce $1$ nie jest zerowa.
Dla dostatecznie małego $\varepsilon$ otrzymujemy
zatem
$\langle T_{\rho_\varepsilon^2},\varphi\rangle
\ge c/(2\varepsilon)\to+\infty$.
Kwadraty tej aproksymacji delty nie mają więc
granicy w $\mathcal D'(\R)$, ponieważ zbieżność
dystrybucyjna wymaga skończonej granicy działania
na każdym ustalonym teście.
Jest to przeszkoda dla definiowania iloczynu
przez zwykłe przejście do granicy gładkich
aproksymacji, odrębna od powyższego argumentu
algebraicznego.
\end{przyklad}

\section[Podsumowanie: przygotowanie do analizy Fouriera]{Podsumowanie i przygotowanie do analizy Fouriera na $\R^d$}

Przestrzeń funkcji testowych ma topologię granicy
induktywnej, której zbieżność ciągów wymaga wspólnego
zwartego nośnika i jednostajnej zbieżności wszystkich
pochodnych (sekcja~\ref{sec:pde7-tests}).
Dystrybucje są ciągłymi funkcjonałami liniowymi,
opisanymi równoważnie przez lokalne oszacowania
~\eqref{eq:pde7-estimate}. Mają pochodne dowolnego
rzędu, ale pochodna słaba funkcji istnieje tylko
wtedy, gdy odpowiednia pochodna dystrybucyjna
jest regularna (sekcja~\ref{sec:pde7-weak}).

Lokalne wygładzanie dystrybucji daje funkcje gładkie
i komutuje z pochodnymi. Dla funkcji z $L^p_{\mathrm{loc}}$,
$1\le p<\infty$, przybliżone jedynki dają również
zbieżność w normie $L^p$ na zwartych podzbiorach,
natomiast w $L^\infty$ taka zbieżność nie zachodzi
na ogół (sekcja~\ref{sec:pde7-mollifiers}).

Rozwiązania fundamentalne operatorów o stałych
współczynnikach pozwalają przez splot z gładkim
wymuszeniem o zwartym nośniku konstruować rozwiązania
na całej przestrzeni (sekcja~\ref{sec:pde7-fundamental}).
Nie zapewnia to samo przez się jednoznaczności
ani spełnienia zadanych warunków brzegowych.
Lemat Weyla pokazuje, że równanie $\Delta T=0$
wymusza gładką i rzeczywiście analityczną
reprezentację dystrybucji
(twierdzenie~\ref{thm:pde7-weyl}).
Nie każdy iloczyn dystrybucji jest jednak określony;
w równaniach nieliniowych należy precyzyjnie
ustalić sens użytych wyrażeń.

W następnym rozdziale~\ref{chap:pde-fourier-Rd}
zbadamy transformatę Fouriera na przestrzeni Schwartza
i dystrybucjach temperowanych. Pozwoli to opisywać
symbole operatorów oraz wybrane jądra potencjałowe
w języku mnożników.
